\documentclass{article}
\usepackage[spanish]{babel}
\usepackage[a4paper, total={6.5 in, 10 in}]{geometry}
\decimalpoint
\usepackage{comment} %Comentarios
\usepackage{listings} %Codigo
\usepackage[dvipsnames]{xcolor} %Paquete de colores en Latex
\usepackage{amsmath} 
\usepackage{amsthm}
\usepackage{amssymb}
\usepackage{amsfonts} %Simbolos exclusivos, como los simbolos de campos
\usepackage{tikz} %Sirve para dibujar muchas cosas en general
%\usetikzlibrary{graphs,graphdrawing} %paquetería que tal vez no es necesaria para poder dibujar gráficas
%\usetikzlibrary{positioning} %Definir nodos y distancias más precisas
\usepackage{blindtext} %Parrafos de texto
\usepackage{multicol} %Columnas
\usepackage[hidelinks]{hyperref} % Carga el paquete base
\hypersetup{
    colorlinks=false,           % Si es false, pone cuadros alrededor de los links
    citebordercolor={0 0 1},    % Color azul para el borde de las citas (RGB decimal)
    linkbordercolor={0 0 1},    % Color azul para los links internos
    pdfborderstyle={/S/S/W 1}   % Estilo del borde (S=Sólido, W=Ancho de 1pt)
}

\usetikzlibrary{shapes.geometric}


\title{Algunos resultados de K-tupla dominación}
\author{Moisés S. S.}
\date{\today}



\begin{document}
	%Comenzamos "formalmente" la escritura del texto desde el principio.
	\maketitle
	\section*{Introduction}

	%investigar como hacer citas correctamente
Se analizan y se muestran resultados sobre $k$-tupla dominación provenientes del artículo \cite{cabrera2022some} y de las referencias de los mismos. Los resultados posteriores buscan mostrar nuevas o mejorar viejas cotas ante el numero de $k$-tupla dominación usando diferentes parametros y cuantificadores para dichas cotas. También se busca la profundización y explicación explicita de las  demostraciones presentadas en el artículo, así como aclaraciones de convenciones asumidas en las referencias y la adaptación de las mismas a la dominación utilizada en el artículo.

	%Se necesita hacer la escritura de cada definición necesaria y agregar todas las que se necesiten o eliminar las redundantes.
	\subsection{Algunos resultados básicos necesarios}

\textbf{Def. 1} Def. de gráfica \newline
\textbf{Def. 2} Def. de vecindades cerradas y abiertas\newline
\textbf{Def. 3} Def. de grados abiertos y grados cerrados.\newline
	%Añadiremos la notación de grado de un vértice sobre un conjunto o en caso de ser necesario, se anexará como definición
\textbf{Def. 4} Def. de grado máximo y mínimo.\newline
\textbf{Def. 5} Def. de subgráfica.\newline
\textbf{Def. 6} Def. de grado de subgráfica respecto a la original.\newline
\textbf{Def. 7} Def. de Restar vértices a una gráfica.\newline
\textbf{Def. 8} Gráfica libre de garras.\newline
\textbf{Def. 9} Definición de (A, B)-aristas.\newline
\textbf{Def. 10} Def. de dominación.\newline
\textbf{Def. 11} Def. de número de dominación.\newline
\textcolor{Violet}{\textbf{Definición 12}} $S$ es un conjunto doble dominante de $G$ si, para cada vértice de $G$ es dominado por al menos dos vértices de $S$ contandose a si mismo.\\

%Numero de doble dominación /Usa vecindad uy grado cerrado
\textcolor{Violet}{\textbf{Definición 13}} El número de doble dominación, está definido por la cardinalidad mínima del conuunto doble dominante y es denotado por $\gamma_{\times k}(G)$\\

\textbf{Def. 14} Def. de $k$-dominación.\newline
\textbf{Def. 15} Def. de número de $k$-dominación.\newline
\textbf{Def. 16} Def. de $l$-total $k$- dominación.\newline
	%K-tupla dominación es un caso particular de la definición anterior
\textbf{Def. 17} Def. de $k$-tupla dominación.\newline
	%Hacemos nota e incapie sobre que valorees puede tomar k para que las definiciones sean equivalentes.
\textbf{Def. 18} Def. numero de $k$-tupla dominación.\newline

	$$FinDeDefiniciones$$

%Teorema de dominanción
\textbf{Teo} 
Si $G$ es una gráfica, entonces

$$\frac{n}{1+\Delta(G)}\leq \gamma(G)\leq n-\Delta(G)$$
%Teorema y definiciones más especificas
\textbf{Teo} 

	%Comenzamos la escritura de los primeros resultados y su uso, además de hacerlo en un solo sentido

\newpage

\section{Resultados relacionados con otros numeros de dominación.}

Comenzamos el estudio del número de $K$-tupla dominación respecto a cotas superiores. Para este capitulo, nos enfocaremos en buscar nuevas cotas superiores en términos de dos enteros $k$ y $k'$, además de utilizar los numeros de dominación y $k$-dominación para mejorar nuestro acotamiento, así como mostrar ejemplos puntuales donde las cotas son justas para alguna gráfica particular o alguna familia de gráficas.\\

	%Falta revisar a profundidad este resultado y su demostración.
La demostración del siguiente teorema se hace en una publicación anterior del mismo autor. Este resultado era un problema abierto y recopilado en el año 2020 \cite{hansberg2020multiple}. Este resultado principalmente fue planteado para $k=2$ Posteriormente se demostró dos años depués para cualquier $k\geq 2$, además de dar una mejora considerable a la cota original.\\

\textbf{Problema} Dar una cota superior para $\gamma_{\times k}(G)$ en términos de $\gamma_k(G)$ para cualquier gráfica $G$ con $\delta(G)\geq k-1$.\\

El problema tiene una solución relativamente simple, dada la relación de la $k$-dominación y la $k$-tupla dominación, tenemos la siguiente desigualdad:\\

$$\gamma_{\times k}(G)\leq k\gamma_k(G)$$

Probemos que esta desigualdad se cumple.\\
Sea $D$ un $\gamma_k$-conjunto y definimos al conjunto $D'\subseteq V(G)$ tal que $D\subseteq D'$ y $\delta_{D'}(x) \geq k-1$ para toda $x\in D$. Como $D$ es un conjunto $k$-dominante, $\forall x \in (V(G)\setminus D)$, $\delta_{D'}(x)\geq k$ y por la definición del conjunto $D'$, $\forall v\in D$, $\delta_{D'}(v)\geq k-1$, por lo que, $D'$ es un conjunto $k$-tupla dominante de $G$.\\


Por último, por la construcción del conjunto $D'$, a lo más, se agregaron $k-1$ vértices más al conjunto $D$ para cada vértice del mismo, es decir, a lo más se agregaron $(k-1)|D|$ vértices, por lo que, $|D'|=|D|+|D'\setminus D|$. Entonces:

\begin{align*}
\gamma_{\times k}(G)&\leq |D'|\\
&\leq |D|+|D'\setminus D|\\
&\leq |D|+(k-1)|D|\\
\Longrightarrow \gamma_{\times k}(G)&\leq k\gamma_k(G)
\end{align*}


Ya que se resolvió el problema anterior y se otorgó una justificación a la misma, posterioremente podemos hacer un refinamiento de la desigualdad anteior, haciendolo mejorar en al menos 1. El refinamiento mencionado se cita y se demuestra en el siguiente teorema \cite{cabrera2022note}.\\

\textbf{Teorema 1} Sea $k \geq 2$ entero. Para cualquier gráfica $G$ con $\delta(G) \geq k-1$.

$$\gamma_{\times k}(G) \leq k\gamma_{k}(G)-(k-1)^2$$

\textbf{Dem.}\\

Sea $D$ un $\gamma_k(G)$-conjunto. Sabemos que $\gamma_{\times k}(G)\leq |V(G)|$ y además, tenemos que considerar dos casos. El primer caso es cuando $k|D|-(k-1)^2\ > |V(G)|$, lo que termina la prueba en este caso. El segundo caso  es cuando $k|D|-(k-1)^2\leq |V(G)|$, que será el caso qeu supondremos para el resto de la demostración.\\

Como $D$ es un conjunto $k$-dominante, sabemos que $|D|\geq k$, de lo que podemos deducir los iguiente:\\

%Se demuestra la existencia del conjunto U
\begin{align*}
|D|&\geq k\\
\Longrightarrow |D|-1 &\geq k-1\\
\Longrightarrow (k-1)(|D|-1) &\geq (k-1)^2\\
\Longrightarrow (k-1)|D|-(k-1) &\geq (k-1)^2\\
\Longrightarrow (k-1)|D|-(k-1)^2 &\geq k-1\\
\Longrightarrow k|D|-(k-1)^2 -|D|&\geq k-1\\
\Longrightarrow V(G)-|D|&\geq k|D|-(k-1)^2 -|D|\\
\Longrightarrow V(G)-|D|&\geq k-1
\end{align*}
\newline
Ahora, como vimos anteriormente, podemos encontrar al menos $k-1$ vértices fuera del conjunto, por lo que podemos definir los siguientes conjuntos. Sean $U=\{u_1, ..., u_{k-1}\}\subset V(G)\setminus D$, $D'=D\cup U$ y $D_0=\{v\in D : \delta_{D'}(v)<k-1\}$. \\
Una vez definidos los conjuntos anteriores, podemos deducir las siguientes desigualdades, las cuales provienen del conteo de aristas, la primera entre los conjuntos $U$ y $D_0$ y la segunda desigualdad de los conjuntos $U$ y $D\setminus D_0$.

\begin{figure}

\begin{center}
\begin{tikzpicture}[scale=1.2]
    % --- DIAGRAMA SUPERIOR ---
    % Elipses principales D y U
    \draw (0,0) ellipse (1cm and 1.8cm) node[above=2.1cm] {$D$};
    \draw (3.5,0) ellipse (0.8cm and 1.8cm) node[above=2.1cm] {$U$};

    %Nodos en U
    \filldraw (3.3,0.85) circle (2pt) node (p8) {};
    \filldraw (3.6,0.85) node (p81) {$u_1$};

    %Nodos de densidad
    %\filldraw (3.3,0.25) circle (0.3pt) node (p9) {};
    %\filldraw (3.3,0.05) circle (0.3pt) node (p10) {};
    %\filldraw (3.3,-0.15) circle (0.3pt) node (p11) {};


    \filldraw (3.3,-0.65) circle (2pt) node (p12) [above=0.55cm]{$\vdots$};
    \filldraw (3.7,-0.65) node (p121) {$u_{k-1}$};


    
    %	Nodos en D
    \filldraw (0.2,1.2) circle (2pt) node (p7) {};

    %Nodos de densidad (pequeños)
   % \filldraw (0,1.2) circle (0.3pt) node (p6) {};
    %\filldraw (-0.15,1.2) circle (0.3pt) node (p5) {};
    %\filldraw (-0.3,1.2) circle (0.3pt) node (p4) {};
    
    \filldraw (-0.5,1.2) circle (2pt) node (p3) [right]{$\; \dots$};


    % Subconjunto D_0 dentro de D
    \draw (0,-0.6) ellipse (0.6cm and 1cm) node[above=0.5cm, left=0.1cm] {$D_0$};
    
    % Punto dentro de D_0
    \filldraw (0.2,-0.5) circle (2pt) node (p1) {};
    
    % aristas azules (D_0 - U)
    \draw[Aquamarine] (p1) -- (3.2, 0.8);
    \draw[Aquamarine] (p1) -- (3.2, 0);
    \draw[Aquamarine] (p1) -- (3.2, -0.6);

    % aristas azules (D - D_0)
    \draw[BurntOrange] (p1) -- (0.2, 1);
    \draw[BurntOrange] (p1) -- (0, 1);
    \draw[BurntOrange] (p1) -- (-0.2, 1);
    \draw[BurntOrange] (p1) -- (-0.4, 1);

    % aristas verdes (D - U)
    \draw[LimeGreen] (p7) -- (3.2, 0.8);
    \draw[LimeGreen] (p7) -- (3.2, 0);
    \draw[LimeGreen] (p7) -- (3.2, -0.6);
    

    % --- SEPARACIÓN ---
   % \begin{scope}[yshift=-5.5cm]
    
    % --- DIAGRAMA INFERIOR ---
    % Elipses principales X e Y
    %\draw (0,0) ellipse (1cm and 1.8cm) node[above=1.8cm] {$X$};
    %\draw (3.5,0) ellipse (1cm and 1.8cm) node[above=1.8cm] {$Y$};
    
    % Subconjunto N(S) dentro de Y
    %\draw (3.5,-0.5) ellipse (0.6cm and 0.8cm) node[above=0.4cm] {$N(S)$};
    
    % Punto dentro de N(S)
    %\filldraw (3.3,-0.7) circle (2pt) node (p2) {};
    
    % Líneas azules (aristas)
    %\draw[blue!70!black] (p2) -- (0.5, 0.6);
    %\draw[blue!70!black] (p2) -- (0.7, -0.1);
    %\draw[blue!70!black] (p2) -- (0.5, -0.8);
    
    %\end{scope}
\end{tikzpicture}
\end{center}

\caption{Conjuntos de vértices definidos en la gráfica $G$ y las posibles aristas entre los mismos}
    \label{fig:grafica_1}
\end{figure}


$$\sum_{v\in D_0}\delta_{D'}(v)\geq \sum_{i=1}^{k-1}\delta_{D_0}(u_i) \hspace{2.5cm} (1)$$
%\newline

El lado izquierdo de la desigualdad, son todas las aristas posibles entre los vértices del conjunto $D_0$ y el conjunto $D'$, apoyandonos de la Figura 1, son las aristas verdes y las aristas azules. El lado derecho de la desigualdad, son todas las posibles aristas entre los vértices de $U$ y los vértices de $D_0$, una vez más, apoyandonos de la Figura 1, son las aristas de color azul. Además, sabemos que $U\subseteq D'$, lo que justifica la desigualdad.

$$|D\setminus D_0|(k-1)\geq \sum_{i=1}^{k-1}\delta_{D\setminus D_0}(u_i) \hspace{2cm} (2)$$


Sabemos que $\forall x \in D\setminus D_0$, tiene $\delta_{D'}(x)\geq k$, por lo que, cada vértice en $D\setminus D'$ puede ser tener como vecinos a todos los vértices del conjunto $U$, es decir, pueden haber a lo más $|D\setminus D_0|(k-1)$ aristas posibles. Lo anterior representa el lado izquierdo de la desigualdad y en la Figura 1, son las aristas de color verde. El lado derecho de la desigualdad son la suma de los grados de todos los vértices del conjunto $U$ sobre el conjunto $D\setminus D_0$, el cual queda acotado por $|D\setminus D_0|(k-1)$.\\

Sumando las desiguadades $(1)$ y $(2)$, tenemos que:\\

\begin{align*}
\sum_{v\in D_0}\delta_{D'}(v)+|D\setminus D_0|(k-1)&\geq \sum_{i=1}^{k-1}\delta_{D_0}(u_i)+\sum_{i=1}^{k-1}\delta_{D\setminus D_0}(u_i)\\
& \geq \sum_{i=1}^{k-1}\delta_{D}(u_i)\\
\Longrightarrow \sum_{v\in D_0}\delta_{D'}(v)+|D\setminus D_0|(k-1)&\geq k(k-1)\\
\end{align*}

Como el conjunto $D$ ya es $k$-dominante de $G$, para cualquier vértice $u_i\in U$ debe tener al menos $k$ vecinos dentro de $D$ y como hay $k-1$ vértices en $U$, la suma total de aristas entre $D$ y $U$ debe de ser al menos $k(k-1)$.\\


Despejamos a $\sum_{v\in D_0}\delta_{D'}(v)$ de la desigualdad

\begin{align*}
\Longrightarrow \sum_{v\in D_0}\delta_{D'}(v)\geq & - (|D\setminus D_0|(k-1)) + k(k-1)\\
\text{Multiplicamos por (-) e invertimos la desigualdad}\\
\Longrightarrow -\sum_{v\in D_0}\delta_{D'}(v)\leq & |D\setminus D_0|(k-1)-k(k-1) \hspace{1cm} (3)
\end{align*}
%\newline


Ahora, definimos al conjunto $D''\subseteq V(G)$ como un conjunto de cardinalidad mínima entre todos los super conjuntos $W$ de $D'$ tales que $\delta_W(x)\geq k-1$ para todo $x \in D$. Con la definición del conjunto $D''$, es fácil ver que, para todo vértice $x$ dentro del conjunto $D'\setminus D_0$ tiene $\delta_{D'}(x)\geq k-1$, por lo que, $x \in D''$. Solo resta ver que sucede con lo vértices en $D_0$.\\



\begin{center}
\begin{tikzpicture}[scale=1.2]
    % --- DIAGRAMA SUPERIOR ---
    % Elipses principales D y U
    \draw[rotate=90] (0.1,-1.9) ellipse (2.4cm and 3.4cm) node[above=3cm] {$D''$};
    \draw (0,0) ellipse (1cm and 1.8cm) node[above=2.1cm] {$D$};
    \draw (3.5,0) ellipse (0.8cm and 1.8cm) node[above=2.1cm] {$U$};

    %Nodos en U
    \filldraw (3.3,0.85) circle (2pt) node (p8) {};
    \filldraw (3.6,0.85) node (p81) {$u_1$};

    %Nodos de densidad
    %\filldraw (3.3,0.25) circle (0.3pt) node (p9) {};
    %\filldraw (3.3,0.05) circle (0.3pt) node (p10) {};
    %\filldraw (3.3,-0.15) circle (0.3pt) node (p11) {};


    \filldraw (3.3,-0.65) circle (2pt) node (p12) [above=0.55cm]{$\vdots$};
    \filldraw (3.7,-0.65) node (p121) {$u_{k-1}$};


    
    %	Nodos en D
    \filldraw (0.2,1.2) circle (2pt) node (p7) {};

    %Nodos de densidad (pequeños)
   % \filldraw (0,1.2) circle (0.3pt) node (p6) {};
    %\filldraw (-0.15,1.2) circle (0.3pt) node (p5) {};
    %\filldraw (-0.3,1.2) circle (0.3pt) node (p4) {};
    
    \filldraw (-0.5,1.2) circle (2pt) node (p3) [right]{$\; \dots$};


    % Subconjunto D_0 dentro de D
    \draw (0,-0.6) ellipse (0.6cm and 1cm) node[above=0.5cm, left=0.1cm] {$D_0$};
    
    % Punto dentro de D_0
    \filldraw (0.2,-0.5) circle (2pt) node (p1) {};

    %aristas moradas (D_0)
    \draw[Mulberry] (p1) -- (2,-1.7);
    \draw[Mulberry] (p1) -- (2,-1.3);
    \draw[Mulberry] (p1) -- (2,-2);
    
    % aristas azules (D_0 - U)
    \draw[Aquamarine] (p1) -- (3.2, 0.8);
    \draw[Aquamarine] (p1) -- (3.2, 0);
    \draw[Aquamarine] (p1) -- (3.2, -0.6);

    % aristas azules (D - D_0)
    \draw[BurntOrange] (p1) -- (0.2, 1);
    \draw[BurntOrange] (p1) -- (0, 1);
    \draw[BurntOrange] (p1) -- (-0.2, 1);
    \draw[BurntOrange] (p1) -- (-0.4, 1);

    % aristas verdes (D - U)
    \draw[LimeGreen] (p7) -- (3.2, 0.8);
    \draw[LimeGreen] (p7) -- (3.2, 0);
    \draw[LimeGreen] (p7) -- (3.2, -0.6);
    

    % --- SEPARACIÓN ---
   % \begin{scope}[yshift=-5.5cm]
    
    % --- DIAGRAMA INFERIOR ---
    % Elipses principales X e Y
    %\draw (0,0) ellipse (1cm and 1.8cm) node[above=1.8cm] {$X$};
    %\draw (3.5,0) ellipse (1cm and 1.8cm) node[above=1.8cm] {$Y$};
    
    % Subconjunto N(S) dentro de Y
    %\draw (3.5,-0.5) ellipse (0.6cm and 0.8cm) node[above=0.4cm] {$N(S)$};
    
    % Punto dentro de N(S)
    %\filldraw (3.3,-0.7) circle (2pt) node (p2) {};
    
    % Líneas azules (aristas)
    %\draw[blue!70!black] (p2) -- (0.5, 0.6);
    %\draw[blue!70!black] (p2) -- (0.7, -0.1);
    %\draw[blue!70!black] (p2) -- (0.5, -0.8);
    
    %\end{scope}

\end{tikzpicture}

\vspace{0.5cm}
Figura 2: \small{El conjunto $D''$ proporciona los vecinos necesarios para los vértices de $D_0$}

\end{center}


Sea $v\in D_0$. Sabemos que todos los vértices $v$ en $D$ tienen $\delta_W(v)\geq k-1$, particularmente, si $v\in D_0$, entonces, el vértice $v$ tiene al menos $k-1-\delta_{D'}(v)$ vecinos en $(W\setminus D')\subseteq (W\setminus D)$, como esto es para cada vértice de $D_0$, tenemos que a lo más son: $|D_0|(k-1)-\sum_{v\in D_0}\delta_{D'}(v)$ vecinos totales. \\
De la misma forma, observamos que al contar los vecinos de cada $v\in D_0$, estos están en $D''\setminus D'$, por lo que, al sumar todos estos vecinos, son a lo menos $|D''\setminus D'|$ ya que $D''$ es de cardinalidad mínima.\\

De lo anterior, tenemos la siguiente desigualdad:\\

\begin{align*}
|D''\setminus D'|&\leq |D_0|(k-1)-\sum_{v\in D_0}\delta_{D'}(v)\\
\text{Utilizando la desigualdad (3):}\\
\Longrightarrow|D''\setminus D'|&\leq |D_0|(k-1)+|D\setminus D_0|(k-1)-k(k-1)\\
\Longrightarrow|D''\setminus D'|&\leq (|D_0|+|D\setminus D_0|)(k-1)-k(k-1)\\
\Longrightarrow|D''\setminus D'|&\leq (|D|)(k-1)-k(k-1) \hspace{1.5cm} (4)\\
\end{align*}

Por la definición del conjunto $D''$, para cada $x\in D$, $\delta_{D''}(x)\geq k-1$ y, como $D$ es $\gamma_k(G)$-dominante y $D\subset D''$, para cualquier $y\in V(G)\setminus D$, $\delta_{D''}\geq k$, entonces, esto es equivalente a, para todo vértice $x\in V(G)$, $\delta_{D''}(v)\geq k-1$, por lo tanto, $D''$ es un conjunto $k$-tupla domiante de $G$ y $\gamma_{\times k}(G)\geq |D''|$. \\

Por último, recordemos las cardinalidades de $D''$ $D'$ y posteriormente sus tituiremos estos valores sobre la desigualdad (4), $|D''|=|D'|+|D''\setminus D'|$ y $|D'|=|D|+k-1$, entonces: $|D''\setminus D'|=|D''|-|D'|$\\


\begin{align*}
|D''|-|D'|&\leq (|D|)(k-1)-k(k-1)\\
\Longrightarrow |D''|&\leq |D'|+|D|(k-1)-k(k-1)\\
\Longrightarrow|D''|&\leq|D|+k-1+ |D|(k-1)-k(k-1)\\
\Longrightarrow|D''|&\leq |D|+k-1+k|D|-|D|-k^2+k\\
\Longrightarrow|D''|&\leq k|D|-k^2+2k-1\\
\Longrightarrow|D''|&\leq k|D|-(k-1)^2\\
\end{align*}

Como $D$ es $\gamma_k(G)$-dominante y $D''$ es $k$-tupla dominante, podemos concluir que:\\
$$\gamma_{\times k}(G)\leq k\gamma_k(G)-(k-1)^2 $$  \newline

 $\hspace{15.5cm}\blacksquare$\\
%$\blacksquare$


El siguiente resultado muestra otra cota superior, motivada por el teorema anterior, para el número de $k$-tupla dominación de una gráfica $G$ en términos de $k$-dominación y $k'$-tupla dominación con $1 \le k' < k$.\\

\textbf{Teorema 2} Sea $k, k'$ dos enteros tales que $k > k' \geq 1$,
$$\gamma_{\times k}(G) \leq \gamma_{\times k'}(G) + (k - k')\gamma_{k}(G)$$

\textbf{Dem.}\\
Sea $S$ un $\gamma_{\times k'}(G)$-conjunto y sea $D$ un $\gamma_{k}(G)$-conjunto. Notemos lo siguiente:\\

\begin{center}
$\forall v \in V(G), \delta_{S}[v] \geq k'$ \hspace{0.2cm} \text{y} \hspace{0.2cm}
$\forall u \in V(G)\setminus D, \delta_{D}(u) \geq k$
\end{center}

Sea $D_{0} = \{v \in D : \delta_{D \cup S}(v) < k-1\}$.
	%D_{0} \subset D
Si $v\in D_0\setminus S$, como $S$ es $k'$-tupla dominante, $\delta_S(v)\ge k'$, así $\delta_{S\cup D}(v)\ge k'$. De igual forma, si $v\in D_0\cap S$ tenemos que $\delta_S(v) \ge k'-1$, lo que implica que $\delta_{S\cup D}(v)\ge k'-1$.


De lo anterior, podemos deducir la siguientes desigualdades:

$$\sum_{v \in D_{0}\setminus S} \delta_{D\cup S}(v) \geq k'|D_{0}\setminus S|,  \sum_{v \in D_{0}\cap S} \delta_{D\cup S}(v) \geq (k'-1)|D_{0} \cap S|$$

Sumando ambas desigualdades, tenemos

$$\sum_{v\in D_{0}\setminus S} \delta_{D \cup S}(v) + \sum_{v\in D_{0} \cap S} \delta_{D \cup S}(v) \geq k'|D_{0} \setminus S| + (k'-1)|D_{0} \cap S|$$

\noindent
\begin{align*}
\Longrightarrow \sum_{v\in D_{0}} \delta_{D \cup S}(v) &\geq k'(|D_{0}| - |D_{0} \cap S|) + (k'-1)(|D_{0}| + |S| - |D_{0} \cup S|)\\
\Longrightarrow \sum_{v\in D_{0}} \delta_{D \cup S}(v) &\geq k'|D_{0}| - k'|D_{0} \cap S| + (k'-1)|D_{0}| + (k'-1)|S| - (k'-1)|D_{0} \cup S|\\
\Longrightarrow \sum_{v\in D_{0}} \delta_{D \cup S}(v) &\geq k'|D_{0}| - k'|D_{0}| - k'|S| + k'|D_{0} \cup S| + (k'-1)|D_{0}| + (k'-1)|S| - (k'-1)|D_{0} \cup S|\\
\Longrightarrow \sum_{v\in D_{0}} \delta_{D \cup S}(v) &\geq |D_{0} \cup S| + (k'-1)|D_{0}| - |S|\\
\Longrightarrow \sum_{v\in D_{0}} \delta_{D \cup S}(v) &\geq (k'-1)|D_{0}| + |D_{0}\setminus S| & (2.1)
\end{align*}
Como $|D \setminus D_{0}| \geq |D \setminus (D_{0} \cup S)|$ y por hipotésis $k > k'$, tenemos que:
    
\begin{align*}
    & (k-1)|D \setminus D_{0}| \geq (k'-1)|D\setminus D_{0}| + |D\setminus D_{0}|\geq (k'-1)|D\setminus D_{0}| + |D\setminus (D_{0} \cup S)|  &  (2.2)
\end{align*}
    
De las desigualdades 2.1 y 2.2, sumamos y obtenemos lo siguiente:
    
\begin{align*}
\Longrightarrow & \sum_{v\in D_{0}} \delta_{D \cup S}(v) + (k-1)|D \setminus D_{0}| \geq (k'-1)|D_{0}| + |D_{0}\setminus S| + (k'-1)|D\setminus D_{0}| + |D\setminus (D_{0} \cup S)|\\
\Longrightarrow & \sum_{v\in D_{0}} \delta_{D \cup S}(v) + (k-1)|D \setminus D_{0}| \geq (k'-1)(|D_{0}| + |D\setminus D_{0}|) + |D_{0}\setminus S| + |D \setminus (D_{0}\cup S)|\\
\Longrightarrow & \sum_{v\in D_{0}} \delta_{D \cup S}(v) + (k-1)|D \setminus D_{0}| \geq (k'-1)|D| + |D\setminus S| \\
\Longrightarrow & \sum_{v\in D_{0}} \delta_{D \cup S}(v) \geq (k'-1)|D| + |D\setminus S| - (k-1)|D \setminus D_{0}| & (2.3)
\end{align*}
    %Esto último se cumple ya que por hip. $\delta(G) \geq k-1$, por lo que, podemos considerar todos los vértices que puedan cumplir esta caracteristica.

Denotamos por $W$ a cualquier supra conjunto de $D\cup S$ tal que, $\delta_W(v)\ge k-1$ para todo $v$ en $D$, y sea $W'$ uno de estos $W$ tal que $W'$ es de cardinalidad mínima. Por la hipotésis de $\delta(G)\geq k-1$ y por definición del conjunto $D_{0}$, implica que para todo $x \in D\setminus D_{0}$, $\delta_{D\cup S}(x) \geq k-1$, de esta forma podemos deducir que, para todo $v \in D_{0}$, tienen al menos $k-1-\delta_{D\cup S}(v)$ vecinos en el conjunto $W\setminus (D\cup S)$, por lo que, a lo más pueden haber $|D_{0}|(k-1) - \sum_{v\in D_{0}} \delta_{D\cup S}(v)$ en el conjunto $W\setminus (D\cup S)$, y en particular para $W'\setminus (D\cup S)$, entonces:\\


\begin{center}
\begin{tikzpicture}[scale=1.2]
    % --- DIAGRAMA SUPERIOR ---
    % Elipses principales D y U
    \draw[rotate=90] (0.1,-1.9) ellipse (2.4cm and 3.4cm) node[above=3cm] {$D''$};
    \draw (0,0) ellipse (1cm and 1.8cm) node[above=2.1cm] {$D$};
    \draw[rotate=90] (-1,-2.3) ellipse (1cm and 2cm) node[above=1.3cm, right=1.3cm] {$S$};

    %Nodos en U
    \filldraw (1.5,-1.4) circle (2pt) node (p8) {};
    \filldraw (1.5,-1.7) node (p81) {$u_1$};

    %Nodos de densidad
    %\filldraw (3.3,0.25) circle (0.3pt) node (p9) {};
    %\filldraw (3.3,0.05) circle (0.3pt) node (p10) {};
    %\filldraw (3.3,-0.15) circle (0.3pt) node (p11) {};


    \filldraw (3.2,-1.4) circle (2pt) node[rotate=90] (p12) [above=0.6cm]{$\vdots$};
    \filldraw (3.2,-1.7) node (p121) {$u_{k-1}$};


    
    %	Nodos en D
    \filldraw (0.2,1.2) circle (2pt) node (p7) {};

    %Nodos de densidad (pequeños)
   % \filldraw (0,1.2) circle (0.3pt) node (p6) {};
    %\filldraw (-0.15,1.2) circle (0.3pt) node (p5) {};
    %\filldraw (-0.3,1.2) circle (0.3pt) node (p4) {};
    
    \filldraw (-0.5,1.2) circle (2pt) node (p3) [right]{$\; \dots$};


    % Subconjunto D_0 dentro de D
    \draw (0,-0.6) ellipse (0.6cm and 1cm) node[above=0.5cm, left=0.1cm] {$D_0$};
    
    % Punto dentro de D_0
    \filldraw (0,-0.2) circle (2pt) node (p1) {};

    %aristas moradas (D_0)
    \draw[Mulberry] (p1) -- (2,1.3);
    \draw[Mulberry] (p1) -- (2,1);
    \draw[Mulberry] (p1) -- (2,1.7);
    
    % aristas azules (D_0 - S)
    \draw[Aquamarine] (p1) -- (1.45,-1.4);
    \draw[Aquamarine] (p1) -- (3.2,-1.35);
    \draw[Aquamarine] (p1) -- (2.5,-1.35);

    % aristas naranjas (D - D_0)
    \draw[BurntOrange] (p1) -- (0.2, 1);
    \draw[BurntOrange] (p1) -- (0, 1);
    \draw[BurntOrange] (p1) -- (-0.2, 1);
    \draw[BurntOrange] (p1) -- (-0.4, 1);

    % aristas verdes (D - S)
    \draw[LimeGreen] (p7) -- (1.45,-1.4);
    \draw[LimeGreen] (p7) -- (3.2,-1.35);
    \draw[LimeGreen] (p7) -- (2.5,-1.35);
    

    % --- SEPARACIÓN ---
   % \begin{scope}[yshift=-5.5cm]
    
    % --- DIAGRAMA INFERIOR ---
    % Elipses principales X e Y
    %\draw (0,0) ellipse (1cm and 1.8cm) node[above=1.8cm] {$X$};
    %\draw (3.5,0) ellipse (1cm and 1.8cm) node[above=1.8cm] {$Y$};
    
    % Subconjunto N(S) dentro de Y
    %\draw (3.5,-0.5) ellipse (0.6cm and 0.8cm) node[above=0.4cm] {$N(S)$};
    
    % Punto dentro de N(S)
    %\filldraw (3.3,-0.7) circle (2pt) node (p2) {};
    
    % Líneas azules (aristas)
    %\draw[blue!70!black] (p2) -- (0.5, 0.6);
    %\draw[blue!70!black] (p2) -- (0.7, -0.1);
    %\draw[blue!70!black] (p2) -- (0.5, -0.8);
    
    %\end{scope}

\end{tikzpicture}

\vspace{0.5cm}
Figura 3: \small{El conjunto $D''$ proporciona los vecinos necesarios para los vértices de $D_0$}.

\end{center}


\begin{align*}
|W'\setminus (D\cup S)| & \leq |D_{0}|(k-1) - \sum_{v\in D_{0}} \delta_{D\cup S}(v)\\
\end{align*}

Utilizando la desigualdad (2.3) tenemos: 

\begin{align*}
	|W'\setminus (D\cup S)| & \leq (k-1)|D_{0}| - ((k'-1)|D| + |D\setminus S| - (k-1)|D \setminus D_{0}|) \\
    & \leq (k-1)|D_{0}| - (k'-1)|D| - |D\setminus S| + (k-1)|D\setminus D_{0}|\\
    & \le (k-1)|D| - (k'-1)|D| - |D\setminus S|\\
    & \le (k - k')|D| - |D \setminus S|
\end{align*}

    %Se verifica que W es conjunto k-tupla dominante
Ahora veamos que $W'$ es un conjunto $k$-tupla dominante. \\
Por como se definió $W'$, si $v\in D \subseteq W'$, tenemos que $\delta_{W'}(v) \geq k-1$, por lo que $\delta_{W'}[v]\geq k$. Si $v\notin D$, como $D\subset W'$ y $D$ es $k$-dominante, entonces, $\delta_{W}(v) \geq k$.

Por lo que $W'$ es un conjunto $k$-tupla dominante. Por lo tanto:
    
\begin{align*}
    %Desigualdad final 
    \gamma_{\times k}(G) &\leq |W'| \\
	&\le |S|+ |D\setminus S| + |W'\setminus (D\cup S)|\\
    & \le |S|+ |D\setminus S| + ((k-k')|D|-|D\setminus S|)\\
    & \le |S| + (k-k')|D|\\
    & \le \gamma_{\times k'}(G) + (k-k')\gamma_{k}(G)
\end{align*}
    $\hspace{15.5cm}\blacksquare$

Podemos notar que la cota superior dad en el teorema 2 es justa para algunas gráficas con una $k$ y $k'$ en particular. Por ejemplo, consideramos $k'=1$ y $k=k'+1=2$, para la siguiente gráfica:\tikz \draw[rotate=30][white] (0,0) ellipse [x radius=6pt, y radius=3pt];\\



%Dibujo de la gráfica, dos ruedas de orden 5 con los centros conectados.
\begin{center}
\begin{tikzpicture}[scale=2.2,
    every node/.style={circle, draw, fill=white, inner sep=1.5pt}]

% --- Lado izquierdo, Rueda de orden 5 ---
\node (a1)[fill=blue!20] at (0,1) {$v_1$};
\node (a2)[fill=blue!20] at (1,1) {$v_2$};
\node (a3)[fill=blue!20] at (1,0) {$v_3$};
\node (a4)[fill=blue!20] at (0,0) {$v_4$};
\node (c1)[fill=orange!40] at (0.5,0.5) {$x$};

\draw (a1)--(a2)--(a3)--(a4)--(a1);
\draw (c1)--(a1);
\draw (c1)--(a2);
\draw (c1)--(a3);
\draw (c1)--(a4);

% --- lado derecho Rueda de orden 5 ---
\node (b1)[fill=green!20] at (3,1) {$u_1$};
\node (b2)[fill=green!20] at (4,1) {$u_2$};
\node (b3)[fill=green!20] at (4,0) {$u_3$};
\node (b4)[fill=green!20] at (3,0) {$u_4$};
\node (c2)[fill=orange!40] at (3.5,0.5) {$y$};

\draw (b1)--(b2)--(b3)--(b4)--(b1);
\draw (c2)--(b1);
\draw (c2)--(b2);
\draw (c2)--(b3);
\draw (c2)--(b4);

% --- Arista que conecta los centros ---
\draw (c1)--(c2);

\end{tikzpicture}


\vspace{0.5cm}
Figura 4: \small{Gráfica $G$ para la que se cumple la igualdad del teorema 2, considerando $k=2$}.


\end{center}
%\newpage
	
	%Ejemplo directo del teorema 2.2 para el que la desigualdad es exacta

Para la gráfica en la figura 4, sean $X=\{x,y\}$ conjunto dominante mínimo, $W=\{v_2, v_4, u_4, u_2\}$ conjunto $2$-dominante mínimo y $T=\{v_2, v_3, x, u_1, u_4, y\}$ conjunto $2$-tupla dominante mínimo. Por lo que: $\gamma(G)=2$, $\gamma_{2}(G)=4$ y $\gamma_{\times 2}(G)=6$, entonces:
	
	
	$$\gamma(G) + (2-1)\gamma_{2}(G) = 2+4 = 6 = \gamma_{\times 2}(G)$$
	

Como consecuencia directa del teorema 2, tenemos el siguiente resultado para $k'=1$, el cual presenta una mejora ante el resultado obtenido en el teorema 1, siempre que $\gamma(G)<\gamma_{k}-(k-1)^2$, lo que es equivalente a que $\gamma(G)+(k-1)^2<\gamma_{k}(G)$.\\

	%Teo 2.2 para el caso K'=1
\textbf{Corolario 2.1} Sea $k\geq 2$ un entero. Para cualquier gráfica $G$ con $\delta(G)\ge k-1$

$$\gamma_{\times k}(G) \leq (k-1)\gamma_{k}(G) + \gamma(G).$$

Continuando con otro corolario derivado del teorema 2, el cual considera el caso $k'=k-1$.\\
	%consultar las referencias citadas y los resultados previamente mencionados.
	
	%teorema 2.2 para el caso k'=k-1
\textbf{Corolario 2.2} Sea $K \geq 2$ un entero. Para cualquier gráfica $G$ con $\delta(G) \geq k-1$ \\

$$\gamma_{\times k}(G) \leq \gamma_{\times (k-1)}(G)+\gamma_{k}(G).$$ 
	
	
\begin{comment}
	Cambiamos la convención utilizada en el artículo original.
	$G_{k,r}=k_{k-1} \lor N_{r}$ \\\\
	$N_{r}$ va a ser la gráfica $(K_{r})^{c}$, la gráfica de orden $r$ sin aristas.
\end{comment}
	
Veamos que la desigualdad anterior es justa. Para cualesquiera $k,r \in \mathbb{Z}$, tal que $r\geq k \geq 2$, sea $G_{k,r}$ la gráfica definida por la suma de $K_{k-1}$ y $(K_{r})^{c}$, es decir, $G_{k,r}=K_{k-1}+(K_{r})^{c}$. Podemos deducir lo siguiente.

\begin{align*}
\gamma_{\times k}(G_{k,r}) & = |V(G_{k,r})|=k-1+r,\\
\gamma_{\times (k-1)}(G_{k,r}) & =k-1,\\
\gamma_{k}(G_{k,r}) & =r\\
%Por \ el \ corolario \ 2.4 \ tenemos \ lo \ siguiente:\\
\Longrightarrow \gamma_{\times k}(G_{k,r})= k-1+r & = k-1+r = \gamma_{\times (k-1)}(G_{k,r})+\gamma_k(G_{k,r})
\end{align*}

% Podemos mostrar un ejemplo más númerico que pueda funcionar.
Veamos un ejemplo particular de esta familia de gráficas. Consideramos $k=3$ y $r=4$, en el que podemos exhibir los conjuntos 3-tupla dominante, 2-tupla dominante y 3-dominante. 

\begin{center}
    \begin{tikzpicture}[scale=1, 
    k3node/.style={circle, draw=red!120!black, fill=red!20, thick, minimum size=6mm},
    k4node/.style={circle, draw=blue!120!black, fill=blue!20, thick, minimum size=6mm}]

    % ==========================================
    % 1. Vértices del conjunto K_3 (Triángulo)
    % ==========================================
    %\node[k3node] (v1) at (0, 2) {$1$};
    \node[k3node] (v2) at (-1, 3.5) {1};
    \node[k3node] (v3) at (1, 3.5) {2};

    % ==========================================
    % 2. Vértices del conjunto K^c_4 (Independiente)
    % ==========================================
    \node[k4node] (u1) at (-2.5, 0) {3};
    \node[k4node] (u2) at (-0.8, 0) {4};
    \node[k4node] (u3) at (0.8, 0) {5};
    \node[k4node] (u4) at (2.5, 0) {6};

    % ==========================================
    % 3. Aristas internas de K_3
    % ==========================================
    \draw[thick, red!100!black] (v2) -- (v3);

    % ==========================================
    % 4. Aristas de la suma (Join)
    % Conectamos cada vértice de K_3 con cada vértice de K^c_4
    % ==========================================
    % Ponemos las aristas del join en gris y un poco más delgadas 
    % para que no saturen visualmente la gráfica
    \foreach \v in {v2, v3} {
        \foreach \u in {u1, u2, u3, u4} {
            \draw[black] (\v) -- (\u);
        }
    }

    \end{tikzpicture}
    
    \vspace{0.5cm}

    Figura 5: \small{Para los enteros $k=3$ y $r=4$, consideramos la siguiente gráfica $G= K^{c}_{4} + K_{2}$, en la que se puede observar a $V(G)$ como conjunto 3-tupla dominante, los vértices rojos como el 2-tupla dominante y los vértices azules como el conjunto 3-dominante.}
\end{center}

Sustituyendo tenemos:
    
$$\gamma_{\times (k-1)}(G)+\gamma_{k}(G) = 4+2=6=\gamma_{\times k}(G).$$

\newpage

	
\subsection{Cotas inferiores en términos de otros números de dominación y \texorpdfstring{$l$}{l}-total \texorpdfstring{$k$}{k}-dominación}

	%El autor utiliza un resultado conveniente, pero no se demuestra ni se especifican las condiciones para las cuales se utiliza el resultado original aquí. El resultado usado es el siguiente: \\

La dominación se ha estudiado ampliamente en la literatura lo que tuvo como consecuencia nuevos tipos de dominación con diferentes parámetros a considerar, como hemos visto anteriormente. Hemos utilizado la $k$-dominación, la cual depende de la cantidad de vecinos dominantes $(k)$ que tenga cada vértice que no es dominante, por otro lado, la $k$-tupla dominación también tiene como parámetro esto último y además pedimos una condición interior para los vértices dominantes, pedimos $k-1$ vecinos dominantes para cada vértice dominante. Favaron et all, generalizaron la idea de la $k$-tupla dominación, como $l$-total $k$-dominación \cite{favaron2001domination}. Enunciamos la $l$-total $k$-dominación en la siguiente definición.\\

%Para cualquier gráfica $G$ con $\delta(G) \geq k > k' \geq 1$ \\
%$$\gamma_{\times k}(G) \geq \gamma_{k'}(G)+k-k'$$
	
	
	%Veamos la desigualdad citada y vamos a "demostrar" el resultado para implementarlo aquí. 
	%Definición necesaria para la desigualdad
%Si \delta(G)<min{l,k}, no tendría sentido, pues no alcanzan los vértices para dominar, sobretodo para l (total dominación) 

\textcolor{blue}{\textbf{Definición 1}} Para dos enteros $l \ge 0$, $k>0$ y $X\subseteq V(G)$, decimos que $X$ es $l$-total $k$-dominante de una gráfica $G=(V,E)$, si para cada vértice $v$ en $X$, $v$ tiene al menos $l$ vecinos en $X$ y para cada vértice $v$ en $V \setminus X$, $v$ tiene al menos $k$ vecinos en $X$. \\

	%Vamos a ejemplificar los casos para cada valor de l, k y el grado mínimo de G, en donde tenemos 6 posibles combinaciones.
\textcolor{blue}{\textbf{Nota.}} Si existe un conjunto $X$, $l$-total $k$-dominante, entonces $\delta(G) \geq min\{l,k\}$, por lo que existe al menos un vértice de grado $\delta(G)$ en $X$ o en $V(G) \setminus X$. Notamos que los índices $l$ y $k$ no se les pide ninguna relación entre ellos. De acuerdo a la relación entre $\delta(G)$, $l$ y $k$, podemos tener los siguientes casos posibles:\\


\textcolor{blue}{\textbf{Caso 1.}} $l \le \delta(G)$.\\
El conjunto $V(G)$ es $l$-total $k$-dominante, ya que todos los vértices de $V(G)$ tienen al menos $l$ vecinos en $V(G)$. Así que existe un conjunto $X$ $l$-total $k$-dominante de cardinalidad mínima.\\


\textcolor{blue}{\textbf{Caso 2.}} $\delta(G) < l$ y $k \le \delta(G)$.\\
No siempre es cierto que exista un conjunto $l$- total dominante, ya que nos faltan vértices para satisfacer la $l$-total dominación. Exhibimos un contraejemplo para la existencia del conjunto $X$.


\begin{center}
\begin{tikzpicture}[
    % Aplicamos tu estilo de nodo de 4.2mm
    every node/.style={
        circle,
        draw,
        inner sep=1pt,
        minimum size=3mm,
        fill=white % Para que las líneas no se vean detrás del nodo
    }
]

	
	%\filldraw (-3,2.5) node (g) {$G:$};
    % 1. Posicionar los 6 nodos en un círculo
    % Usamos coordenadas polares (ángulo:radio)
    \foreach \i in {1,...,6} {
        \node (v\i) at ({(\i-1)*360/6}:1.5cm) {$v_\i$};
    }

    \node (w1) at (3.5,2) {$w_1$};
    \node (w2) at (3.5,0) {$w_2$};
    \node (w3) at (3.5,-2) {$w_3$};

    % 2. Conectar cada nodo con todos los demás
    % El bucle \u recorre del 1 al 5
    % El bucle \v recorre desde el siguiente de \u hasta 6
    \foreach \u in {1,...,5} {
        \pgfmathtruncatemacro{\next}{\u+1}
        \foreach \v in {\next,...,6} {
            \draw (v\u) -- (v\v);
        }
    }

    \foreach \x in {3,4,5} {
    	\pgfmathtruncatemacro{\next}{\x+1}
    	\foreach \v in {\next,...,5} {
    		\draw[red] (v\x) -- (w1);
    	}
    }

    \foreach \x in {1,2,6} {
    	\pgfmathtruncatemacro{\next}{\x+1}
    	\foreach \v in {\next,...,5} {
    		\draw[red] (v\x) -- (w3);
    	}
    }

    \draw[red] (v1) -- (w2);
    \draw[blue] (w1) -- (w2);
    \draw[blue] (w3) -- (w2);

\end{tikzpicture}

\vspace{0.5cm}

Figura 6\vspace{0.5cm}: \small{Consideramos la gráfica $G$ y $k=2$, $\delta(G)=3$, $l=5$.}

\end{center}

 La gráfica $G$ en la figura $6$, no importa que conjunto de vértices se tome, siempre quedan vértices por cumplir las condiciones de la $5$-total dominación, por lo que, $G$ no tiene un conjunto $5$-total $2$-dominante\\


\textcolor{blue}{\textbf{Caso 3.}} $\delta < l$ y $k > \delta(G)$.\\
No siempre podemos encontrar un conjunto $l$-total $k$-dominante, ya que $\delta(G)$ no nos asegura tener la cantidad de vecinos suficiente para dominar dentro o fuera del conjunto $X$. Exhibimos un contraejemplo para el cual no existe el conjunto $X$.\\


\begin{center}
\begin{tikzpicture}[
    % Aplicamos tu estilo de nodo de 4.2mm
    every node/.style={
        circle,
        draw,
        inner sep=1pt,
        minimum size=3mm,
        fill=white % Para que las líneas no se vean detrás del nodo
    }
]

	
	%\filldraw (-3,2.5) node (g) {$G:$};
    % 1. Posicionar los 6 nodos en un círculo
    % Usamos coordenadas polares (ángulo:radio)
    \foreach \i in {1,...,4} {
        \node (v\i) at ({(\i-1)*360/4}:1.5cm) {$v_\i$};
    }

    \node (w1) at (3.5,2) {$w_1$};
    \node (w2) at (3.5,0) {$w_2$};
    \node (w3) at (3.5,-2) {$w_3$};

    % 2. Conectar cada nodo con todos los demás
    % El bucle \u recorre del 1 al 5
    % El bucle \v recorre desde el siguiente de \u hasta 6
    \foreach \u in {1,...,3} {
        \pgfmathtruncatemacro{\next}{\u+1}
        \foreach \v in {\next,...,4} {
            \draw (v\u) -- (v\v);
        }
    }

    \foreach \x in {1,2,3} {
    	\pgfmathtruncatemacro{\next}{\x+1}
    	\foreach \v in {\next,...,3} {
    		\draw[red] (v\x) -- (w1);
    	}
    }

    \foreach \x in {1,3,4} {
    	\pgfmathtruncatemacro{\next}{\x+1}
    	\foreach \v in {\next,...,3} {
    		\draw[red] (v\x) -- (w3);
    	}
    }

    %\draw[red] (v1) -- (w2);
    \draw[blue] (w1) -- (w2);
    \draw[blue] (w3) -- (w2);

\end{tikzpicture}

\vspace{0.5cm}

Figura 7: \small{Consideramos la gráfica $G$ y $k=3$, $\delta(G)=2$, $l=4$.}

\end{center}



	
	%Esta última nota justificará una hipotésis necesaria para la desigualdad antes mencionada.
	
Ya que hemos visto la definición de la $k$-tupla dominación y sus restricciones, ahora veamos la siguiente afirmación y posterioremente adaptaremos el resultado a la $k$-tupla dominación.\\





	% ________________________________________________________

	%			AFIRMACIÓN 1 ( FALSA )

	%_________________________________________________________



\begin{comment}

\textcolor{blue}{\textbf{Afirmación 1}} Sean $k, l, k' ,l'$ enteros positivos, tales que $k\le l+1$ y $k'\le l'+1$. Si existe un $\gamma_{l,k}(G)$-conjunto para la gráfica $G=(V,E)$, entonces, un $\gamma_{l',k'}(G)$-conjunto existe para cada $l' \leq l$, $k' \leq k$, y además 
$$\gamma_{l',k'}(G) \leq \gamma_{l,k}(G) - min\{l-l', k-k'\}.$$\\

\textbf{Dem:}\\

Si $X$ es un $\gamma_{k,l}(G)$-conjunto y si $Y\subset X$ tal que $|Y| = min\{l - l', k - k'\}$, entonces, el conjunto $X\setminus Y$ es un conjunto $l'$-total $k'$-dominante. Sea $v$ en $X\setminus Y$, sabemos que $v$ pertenece a $X$ y $\delta_X(v)\ge l$, entonces, $\delta_{X\setminus Y}(v)\ge l-(l-l')=l'$, por lo que $X\setminus Y$ es $l'$-total dominante.\\

Ahora, sea $u$ en $V\setminus X$, como $X$ es $l$-total $k$-dominante, $\delta_X(u)\ge k$, entonces, $\delta_{x\setminus Y}(u) \ge k-(k-k') = k'$, de esta forma, solo queda el caso de los vértices que están en el conjunto $Y\setminus X$. Sea $w$ en $Y\setminus X$, sabemos que $Y\subset X$, por lo que $\delta_X(w)\ge k$, entonces $\delta_{x\setminus Y}(w)\ge k-|Y|$, Por como se definió el conjunto $Y$, tenemos dos posibles casos.\\

\textbf{Caso 1.} $|Y|= l-l'$, entonces $\delta_{X\setminus Y}(w) \ge l-(l-l')+1 = l'+1$, por hipotésis, $k'\le l'+1$, por lo tanto $\delta_{x\setminus Y}(w)\ge k'$\\

\textbf{Caso 2.} $|Y|= k-k'$, entonces, $\delta_{X\setminus Y}(w) \ge l-(k-k')+1 = l+1-k+k'$, por hipotésis, $k \le l+1$, por lo que $\delta_{X\setminus Y}(w) \ge k'$.\\

Por lo tanto, el conjunto $X\setminus Y$ es un conjunto $l'$-total $k'$-dominante de $G$, por lo tanto,

$$\gamma_{l',k'}(G)\le |X\setminus Y|= \gamma_{l,k}(G) - min\{l-l', k-k'\}$$

    $\hspace{15.5cm}\blacksquare$ \newline


\textbf{Contraejemplo(s) (?)}\\


Adaptando la $l$-total $k$-dominación a la $k$-tupla dominación, tenemos que $l=k-1$. Adaptando el resultado anterior a la $k$-tupla dominación, $l=k-1$ y $l'=0$, entonces, $k\le l+1 = (k-1)+1 = k$ y $k'\le l' +1 = 0+1 = 0$, además, $\gamma_{xk}(G) = \gamma_{k-1,k}(G)$ y $\gamma_{k'}(G) = \gamma_{0,k'}(G)$. \\

\begin{align*}
\text{Sean } k'\le k, \hspace{0.2cm} 0\le k-1, \hspace{0.2cm}\delta(G)\ge k \ge k' \ge 1\\
k\le l+1, \hspace{0.2cm} k'\le l'+1 \hspace{0.2cm} y \hspace{0.2cm} l'\le l, \hspace{0.2cm} k'\le k\\ 
\text{Por la afirmación anterior, sean \vspace{0.5cm}} k'\le l'+1 \hspace{0.2cm} y \hspace{0.2cm} k\le l+1\\
\text{Entonces } k'\le 1 \hspace{0.2cm} y \hspace{0.2cm} k\le k\\
\end{align*}

Sustituyendo tenemos:\\


\begin{align*}
\gamma_{0,k'}(G) \le \gamma_{k-1,k}(G) & - min\{k-1-0, k-k'\}\\
\gamma_{0,k'}(G) \le  \gamma_{k-1,k}(G)& - min\{k-1, k-k'\}\\
\text{Por la adaptación, $k'\le l'+1$, }\\
\Longrightarrow \gamma_{0,k'}(G) \le \gamma_{k-1,k}(G) & - min\{k-1, k-1\}\\
\Longrightarrow \gamma_{0,k'}(G) \le \gamma_{k-1,k}(G) & - (k-1)\\
\Longrightarrow \gamma_{k'}(G) + (k-1) \le \gamma_{\times k}(G) &\\
\end{align*}


\begin{comment} %Demostración original con la condición faltante (Error)

\textcolor{blue}{\textbf{Afirmación 1}} Si existe un $\gamma_{l,k}(G)$-conjunto para la gráfica $G=(V,E)$, entonces, un $\gamma_{l',k'}(G)$-conjunto existe para cada $l' \leq l$, $k' \leq k$, y además 
$$\gamma_{l',k'}(G) \leq \gamma_{l,k}(G) - min\{l-l', k-k'\}.$$\\
	
\textbf{Dem:}\\
	%En la demostración estamos "quitando" los vértices que nos sobran para poder construir un conjunto que cumpla lo que estamos buscando a partir del conjunto \gamma_{l,k}(G).

Si $X$ es un $\gamma_{k,l}(G)$ conjunto y si $Y \subset X$ tal que $ |Y| = min\{l - l', k - k'\}$, tenemos que para cada vértice $v \in X \setminus Y$, : $\delta_{X\setminus Y}(v) \geq l - l + l' = l'$,  por la cardinalidad con la que se tomó $Y$. 
De la misma forma $\delta_{V(G)\setminus(X \setminus Y)}(v) \geq k - k + k' = k'$.
Por lo que $X \setminus Y$ es un $l'$-total $k'$-dominante conjunto de $G$.\\

Ahora, vamos a utilizar la afirmación anterior adaptado a $k$-tupla dominación.\\

	%sustitución necesaria para usar la desgualdad en k-dominación.
Sean $l = k-1$ y $l'= 0$, entonces\\ 	
\begin{align*}
	\gamma_{\times k}(G) &= \gamma_{k-1,k}(G)\\
	\gamma_{k'}(G) &= \gamma_{0,k'}(G)
\end{align*}

Sustituyendo lo anterior en la \textit{afirmación 1}, tenemos lo siguiente:\\	
	%Es un caso particular de la desigualdad original.

Sean $k'\leq k$, $0 \leq k-1$, tal que $\delta(G) \geq k \geq k'\geq 1$,\\
\begin{align*}
	\Longrightarrow\gamma_{0,k'}(G) &\leq \gamma_{k-1,k}(G) - min\{k-1-0, k-k'\}\\
	\Longrightarrow\gamma_{k'}(G) &\leq \gamma_{\times k}(G) - min\{k-1, k-k'\}\\
	\Longrightarrow \gamma_{k'}(G) &\leq \gamma_{\times k}(G) - min\{k-1, k-k'\}\\
	\Longrightarrow \gamma_{k'}(G) &\leq \gamma_{\times k}(G) - (k-k')\\
	\Longrightarrow \gamma_{k'}(G) + k-k'&\leq \gamma_{\times k}(G)
\end{align*}

%\end{comment}

	%Queda comprobar y corroborar si la desigualdad funciona para k'-tupla dominación en lugar de k'-dominación.

Este resultado anterior puede ser mejorado para gráficas conexas que tengan $diam(G) \geq 5$.\\

\end{comment}

% ________________________________________________________

%TEOREMA 2.5, NO ES MEJORA, ES UN RESULTADO INDEPENDIENTE

%_________________________________________________________



\textbf{Teorema 2.5} Sean $k$, $k'$ dos enteros tales que $k > k' \geq 1$. Para cualquier gráfica conexa $G$ con $\delta(G) \geq k$, se cumple la siguiente desigualdad:\\

$$\gamma_{\times k}(G) \geq \gamma_{k'}(G) + (k-k') \Big\lceil \frac{diam(G) + 1}{5} \Big\rceil$$\\
\textbf{Dem:}\\

Sea $D$ un $\gamma_{\times k}(G)$-conjunto y sea $P$ una geódesica de longitud $r$, siendo $r = diam(G)$, denotada por $P = (v_{0}, v_{1}, v_{2}, ..., v_{r})$ y sea $X = \{v_{0}, v_{5}, v_{10}, ..., v_{5 \lfloor \frac{r}{5} \rfloor}\}$ un subconjunto de vértices de $P$. Por como se tomó $X$, podemos deducir que $|X| = \Big\lceil \frac{diam(G) + 1}{5} \Big\rceil$.\\


Ahora, sea $(v_{i}, v_{i+1}, v_{i+2}, v_{i+3}, v_{i+4}, v_{i+5})$ una subtrayectoria de la trayectoria $P$, con $i \in \{0, ..., r - 5\}$ y supongamos que existen dos vértices $u ,w$ tales que $(v_{i}, u), (w, v_{i+5}) \in E(G)$ y a su vez también $(u, v_{i+2}), (v_{i+3}, w)$, de esta forma, tenemos una alternativa de geodésica. No pueden existir la arista $(u, v_{i+3})$ y la arista $(v_{i+2}, w)$, de lo contrario, en ambos casos formamos una trayectoria con menor distancia que $P$, 
%$(v_{i}, u, v_{i+3}, w, v_{i+5})$ o $(v_{i}, u, v_{i+2}, w, v_{i+5})$
pero esto es una contradicción a la elección de $P$. No existen aristas entre $v_{i}$ y $v_{i+2}$, para cualquier $i \in \{0, ..., r-5\}$ porque $P$ es geodésica. Particularmente, para todo $\{x,y\}\in X$, $N[x]\cap N[y]=\emptyset$ más aún, para cada $u_x \in N(x)$ y para cada $u_y\in N(y)$, tenemos que $N[u_x]\cap N[u_y]=\emptyset$ (*).\\

%Argumento similar respecto a N[u_x]\cap N[u_y]=\emptyset


%Veamos por que podemos asegurar la existencia de $D'$ y por qué está bien definido.

Ahora, probaremos que existe $D'\subset N[X] \cap D$ tal que, $|D'| = (k-k')|X|$ y para cada vértice $x \in X$, $\delta_{D'}[x] = k-k'$. Para todo $x \in X$ sabemos que $\delta_{D}[x] \geq k$ y por hipótesis $k > k'$, de esta forma podemos construir los siguientes conjuntos.\\

%Subconjunto de vecindad cerrada de x con la cardinalidad buscada
Sea $S_{x} \subseteq N[x] \cap D$ tal que $|S_{x}| = k'$, entonces, sea $T_{x} = (N[x] \cap D) \setminus S_{x}$, por lo que, $|T_{x}| = |(N[x] \cap D) \setminus S_{x}| \geq k - k'$. Ahora, sea $H_{x} \subseteq T_{x}$ tal que $|H_{x}| = k - k'$. Podemos construir este conjunto para cualquier vértice $x$, entonces, tomamos la unión de estos conjuntos,

$$ \bigcup\limits_{x \in X} H_{x} $$

Por (*) sabemos que esta unión es ajena, así, por construcción del conjunto $H_x$, determinamos la cardinalidad de la unión de estos conjuntos: 

$$|\bigcup\limits_{x \in X} H_{x}| = (k - k')|X|$$

Ahora, como $D$ es un conjunto $k$-tupla dominante, sabemos que, $|N[x]\cap D|\ge k$ y por hipotésis $k>k'$, entonces, $k>k-k'$, por lo que podemos tomar exactamente $k-k'$ vértices de $N[x]\cap D$. Además, $N[x]\cap N[y] = \emptyset$ para cualesquiera $x,y \in X$, por lo tanto $\delta_{D'}[x]=k-k'$, para cada $x \in X$. Por lo que el conjunto $D'$ existe y está bien definido.\\
	%Termina la afirmación de la existencia del conjunto D'

%Demostración de W como un conjunto k-dominante
Sea $W=D\setminus D'$ y veamos que $W$ es un conjunto $k'$-dominante de $G$. Afirmamos lo siguiente, si $v \in V(G)\setminus W$, entonces $\delta_{D'}[v]\le k-k'$. Sabemos que si $v\in X$, entonces $\delta_{D'}[v]= k-k'$. Por otro lado, si $v \in N(X)$, por lo que, $v$ solo puede ser vecino a lo más de todo $N(x)$ para a lo más exactamente un $x \in X$, así tenemos que $\delta_{D'}[v]\le k-k'$. Análogamente, si $v\in V(G) \setminus D'$, $v$ solo puede ser adyacente a todo $N(x)$ para algún $x \in X$ y $\delta_{D'}(v)\le k-k'$.\\







%Dejamos aquí el saltote que me aventé y justifiqué mal.

\begin{comment}

 Sea $v\in V(G)\setminus W$ y como $D'\subseteq N[X]$ y $\delta_{D'}[x]=k-k'$ para todo $x$ en $X$, así podemos afirmar que $\delta_{D'}[v]\le k-k'$. Sabemos que el vértice $v$ puede estar en $X$ o en $N(X)$ o $(V(G)\setminus N[X]\cap D')$, veamos que se cumple la afirmación encada caso. \\


%Importancia de distancia 5 entre vértices de X y el argumento del grado de cada vértice de G sobre D'
Si $v\in X$, sabemos que $\delta_{D'}[x]=k-k'$, por otro lado, si $v\in N(X)$, anteriormente habíamos demostrado que $v$ pertenece a la vecindad abierta de un sólo vértice $x$ y a lo más ser adyacente a todos los vértices de esta vecindad, por lo que, $\delta_{D'}[v]\le k-k'$. Si $V\in V(G)\setminus N[X]\cap D'$, tenemos de forma análoga al caso anterior, $v$ solo puede ser adyacente a lo más a todo $N[x]$ para un sólo vértice $x$, entonces, $\delta_D'(v)\le k-k'$. \\

\end{comment}


Ahora, sea $v \in V(G)\setminus W$, así tenemos dos casos $v \in V(G)\setminus D$ o $v \in D'$:\\

\textbf{Caso 1.} Sea $v \in (V(G)\setminus D)$. Como $D$ es un conjunto $k$-tupla dominante de $G$, tenemos que $\delta_{D}(v) \geq k$ y $\delta_{D'}(v) \leq k-k'$. Por como se define $W$ y por las desigualdades anteriores tenemos lo siguiente.
	
$$\delta_{W}(v) \geq \delta_{D}(v) - \delta_{D'}(v) \Longrightarrow \delta_{W}(v)\geq k-(k-k')=k'$$

\textbf{Caso 2.} Sea $v \in D'$. Como $D'\subseteq D$, tenemos que $\delta_{D}(v) \geq k-1$ y $\delta_{D'}(v) \leq k-k' - 1$. Por lo que, 
$$\delta_{W}(v) \geq \delta_{D}(v)- \delta_{D'}(v) \Longrightarrow \delta_{W}(v) \geq (k-1)-(k-k'-1)=k'$$

De ambos casos podemos concluir que, el conjunto $W$ es $k'$-dominante de $G$. Por lo tanto:

$$\gamma_{k'}(G)\le |W|=|D|-|D'|=|D|-(k-k')|X|=\gamma_{\times k}(G) - (k-k')\Big\lceil \frac{diam(G) + 1}{5} \Big\rceil$$

$\hspace{15.5cm}\blacksquare$\\

Veamos con ejemplo que la desigualdad del teorema anterior es justa. Sea $G = P_{6} \circ K_{3}$.\\

\begin{center}
\begin{tikzpicture}[
    scale=1.15,
    every node/.style={
        circle,
        draw,
        inner sep=1.2pt,
        minimum size=4.2mm
    }
]

% --- Coordinates ---
\foreach \i/\x in {1/0, 2/1.6, 3/3.2, 4/4.8, 5/6.4, 6/8} {
    \node[fill=white] (u\i) at (\x,2.2) {};
    \node[fill=white] (m\i) at (\x+0.4,1.1) {};
    \node[fill=white] (l\i) at (\x,0) {};
}

% --- Black vertices (as in the image) ---
\foreach \i in {2,5} {
    \node[fill=black] at (u\i) {};
}
\foreach \i in {2,5} {
    \node[fill=black] at (m\i) {};
}
\foreach \i in {2,5} {
    \node[fill=black] at (l\i) {};
}

% --- K3 inside each column ---
\foreach \i in {1,...,6} {
    \draw (u\i)--(m\i)--(l\i)--(u\i);
}

% --- Complete connections between consecutive columns ---
\foreach \i in {1,...,5} {
    \foreach \a in {u,m,l} {
        \foreach \b in {u,m,l} {
            \draw (\a\i)--(\b\the\numexpr\i+1\relax);
        }
    }
}

\end{tikzpicture}

\vspace{0.5cm}

Figura 8: \small{Gráfica $G$ está dada por el producto lexicográfico de $P_6 \circ k_3$}

\end{center}



En esta gráfica, vamos a considerar $k=3$ y $k'=2$, con lo que podemos tomar como conjunto $3$-tupla dominante mínimo a los vértices que están oscurecidos y como conjunto $2$-total dominante, podemos considerar a los vértices oscurecidos superiores e inferiores. Por lo que $\gamma_{\times 3}(P_{6} \circ K_{3})=6$ y $\gamma_{2}(P_{6} \circ K_{3})=4$ y además $diam(P_{6} \circ K_{3})=5$ Sustituyendo los valores en la desigualdad del teorema 2.5, tenemos:

\begin{align*}
6&\geq4+(3-2)\Big\lceil\frac{5+1}{5}\Big\rceil\\
6&\geq4+2\\
6&=6
\end{align*}

%"Generalización" del ejemplo anterior
De forma más general podemos pensar en $P_6 \circ k_k$. Usando $k, k'$, tales que $k \geq k' \geq 1$, por lo que tenemos, $\gamma_{\times k}(P_{6} \circ K_{k})=2k, \gamma_{k'}(P_{6} \circ K_{k})=2k'$ y $diam(P_{6} \circ K_{k})=5$.\\


\textcolor{Green}{\textbf{Fun fact.}} El teorema anterior toma como inspiración una afirmación de la $l$-total $k$-dominación \cite{favaron2001domination}, en el que hemos notado que la afirmación es falsa, he aquí el contraejemplo.\\

\textcolor{blue}{\textbf{Afirmación 1}} Si existe un $\gamma_{l,k}(G)$-conjunto para la gráfica $G=(V,E)$, entonces, existe un $\gamma_{l',k'}(G)$-conjunto para cada $l' \leq l$, $k' \leq k$, y además 
$$\gamma_{l',k'}(G) \leq \gamma_{l,k}(G) - min\{l-l', k-k'\}.$$\\

\textbf{Contraejemplo:}\\

Sea $G=C_8$, ciclo de orden 8 y sean $l=2$, $k=4$, $l'=1$ y $k'=3$. Notemos que cada vértice de $C_8$ tienen grado 2, de esta forma el  $\gamma_{2,4}(C_8)$-conjunto es $V(G)$. Para $l'=1$ y $k'=3$, tenemos que el $\gamma_{1,3}(C_8)$-conjunto existe y es $V(C_8)$. Entonces:

\begin{align*}
\gamma_{1,3}(C_8) &\le \gamma_{2,4}(C_8) - min\{2-1, 4-3\}\\
\Longrightarrow 8 &\le 8-1
\end{align*}


Dibujo\\



\begin{center}
\begin{tikzpicture}[
    scale=1.65,
    % Aplicamos tu estilo de nodo de 4.2mm
    every node/.style={
        circle,
        draw,
        inner sep=1pt,
        minimum size=3mm,
        fill=Thistle % Para que las líneas no se vean detrás del nodo
    }
]

	
	%\filldraw (-3,2.5) node (g) {$G:$};
    % 1. Posicionar los 6 nodos en un círculo
    % Usamos coordenadas polares (ángulo:radio)
    \foreach \i in {1,...,8} {
        \node (v\i) at ({(\i-1)*360/8}:1.5cm) {$v_\i$};
    }


    % 2. Conectar cada nodo con todos los demás
    % El bucle \u recorre del 1 al 5
    % El bucle \v recorre desde el siguiente de \u hasta 6
    \foreach \i in {1,...,7} {
        \pgfmathtruncatemacro{\next}{\i+1}
        %\foreach \v in {\next,...,8} {
            \draw (v\i) -- (v\next);
        %}
    }

    \draw (v8) -- (v1);

\end{tikzpicture}

\vspace{0.5cm}

Figura 9: \small{Ciclo de orden 8, siendo $V(G)$ el $\gamma_{2,4}(G)$-conjunto y a su vez el $\gamma_{1,3}(G)$-conjunto. }


\end{center}

De esta forma hemos determinado el contraejemplo como caso particular, pero podemos generalizarlo para cualquier ciclo de orden $n$, y tomando los parámetros $l=2$ y $k'=3$ tales que $\min\{l-l', k-k'\} = 1$.






\newpage

%Sección con muchos resultados, usando parametros de grados y números de independencia.
\section{Cotas  para el número de k-tupla dominación en términos de orden y número de independencia}

%Comenzamos con la cita 15, que habíamos usado anterioremente, que es el caso mas general de la k-tupla dominación. Vamos a escribir e interpretar los resultados usados para nuestro tema

En la literatura existen varios resultados que acotan a número de dominación clásico $\gamma(G)$ con el número de independencia $\alpha(G)$ y el orden. Es por ello que en est sección veremos como podemos utilizar el orden y el número de independencia para acotar superiormente al número de $k$-tupla dominación.\\


\begin{comment}

En principio, la dominación común está relacionada con la independencia de vértices, incluso hay resultados muy conocidos y comunmente utlizados en la dominación teniendo como parámetros al número de dominación o la relación que existe entre los conjuntos independientes y los conjuntos dominantes de una gráfica. En esta sección veremos como podemos utilizar el número de independencia para poder acotar superiormente al número de $k$-tupla dominación.\\
Comenzamos esta sección con un resultado "trivial" de \cite{favaron2001domination} retomando la $l$-total $k$-dominación.\\

\end{comment}
La siguiente afirmación de \cite{favaron2001domination} retoma la $l$-total $k$-dominación para dar una cota superior en términos del orden y el grado mínimo.\\


%Retomamos l-toalk-dominación
\textcolor{blue}{\textbf{Affirmación 2.2}} Sea $G=(V, E)$ una gráfica de orden $n$, si $\delta(G)\geq max\{l,k\}$, entonces
$$l+1\leq \gamma_{l,k}(G)\leq max\{n-\delta(G)+l, n-\delta(G) +k-1\}$$
\textbf{dem:}\\
%Desigualdad inferior izquierda
%Desigualdad inmediata de definición total dominación
Vamos a probar que $l+1\leq \gamma_{l,k}(G)$. Sea $D\subseteq V(G)$, si $D$ es $\gamma_{l,k}(G)$-conjunto, en particular para cada vértice $v \in D$ sabemos que $\delta(G)\ge \max\{l,k\}$, de esta forma $\delta_D[v]\ge l+1$, por lo que es inmediata la primer desigualdad.\\


%Desigualdad superior derecha
%Utilizamos propiedades de máximos y mínimos para las desigualdades del conjunto X
Ahora veamos la otra desigualdad. Sea $X\subseteq V(G)$ tal que $|X|= max \{n-\delta(G)+l, n-\delta(G)+k-1\}$. Por la propiedad de máximos y mínimos $(a-max\{A, B\} = min\{a-A, a-B\})$, tenemos que:\\
\begin{align*}
|V\setminus X|&= n-max\{n-\delta(G)+l, n-\delta(G)+k-1\}\\
&= min\{n-n+\delta(G)-l, n-n+\delta(G)-k+1\}\\
&= min\{\delta(G)-l, \delta(G)-k+1\}\\
\Longrightarrow \delta(G)-l \ge |V\setminus X|  \hspace{1cm}  & \& \hspace{1cm} \delta(G)-k+1 \ge |V\setminus X| 
\end{align*}
Con las desigualdades anteriores, podemos deducir lo siguiente:\\
%Usamos propiedades de grado y grado sobre un conjunto especifico.
\begin{align*}
\forall v \in X \hspace{.5cm} \delta_{X}(v)\geq \delta(G)-|V(G)\setminus X|&\geq \delta(G)-(\delta(G)-l)=l\\
\forall u \in V\setminus X \hspace{.5cm} \delta_{X}(u)\geq \delta(G)-(|V(G)\setminus X|-1)&\geq \delta(G)-(\delta(G)-k+1-1)=k
\end{align*}

De las desigualdades anteriores, hemos demostrado que $X$ es un conjunto $l$-total $k$-dominante, por lo que

$$\gamma_{l,k}(G)\leq |X|$$


$\hspace{15.5cm}\blacksquare$\\

Como consecuencia de la afirmación, se tiene el \textit{teorema 3.1}, proporcionando un acotamiento para el número de $k$-tupla dominación.\\

\textbf{Teorema 3.1} Para cualquier gráfica $G$ con $\delta(G)\geq k-1$
$$\gamma_{\times k}(G)\leq n(G)- \delta(G)+k-1\leq n(G)$$


%Resultado adaptado al teorema 3.1
Basta con tomar $l=k-1$, así $\gamma_{k-1,k}(G) = \gamma_{\times k}(G)$, $max\{n-\delta(G)+k-1, n-\delta(G)+k-1\} = n-\delta(G)+k-1$ y $\delta(G)\geq k$, entonces

$$k\leq \gamma_{l,k}(G)\leq n-\delta(G)+k-1\le n$$

Que es justamente el resultado presentado como el \textit{teorema 3.1}\\

%Referencia a números de independencia
Como caso particular, buscamos acotar el número de $k$-tupla dominación utilizando el número de independencia sobre gráficas libres de $p$-garras. Se proporciona una cota inferior para $\gamma_{\times k}(G)$ \cite{klasing2004hardness}.\\


%Lema original de la referencia 16
\textcolor{LimeGreen}{\textbf{Lema 7.}} Sea $G=(V,E)$ una gráfica $p$-garra libre, con $p\geq 2$ y sea $k\in \mathbb{Z}^+$ tal que $\delta(G)\geq k-1$, entonces:

$$\frac{k\alpha(G)}{(p-1)}\leq |\gamma_{\times k}|$$

\textbf{Dem:}\\

%Demostración directa por construcción de conjuntos dominante e independiente, así como el conteo de aristas entre estos conjuntos.

Sean $D_{k}^*$ un $\gamma_{\times k}(G)$-conjunto de $G$ y $S$ un $\alpha(G)$-conjunto de $G$ y sea $c_{u}=|D_{k}^*\cap N_{G}[u]|$ para todo $u$ en $S$. Como $D_{k}^*$ es $k$-tupla dominante de $G$, entonces, $c_{u}\geq k$ para cada $u\in S$, por lo que

$$\sum_{u\in S}c_u\geq k|S| \hspace{3cm} (1)$$

%Aristas de S a D_k^*
Sea $d_v=|S\cap N_{G}[v]|, \hspace{.3 cm} \forall v\in D_{k}^*$. Como $G$ es $p$-garra libre, tenemos que para todo vértice $v$ en $D_{k}^*$  tienen a lo más $p-1$ vértices independientes en su vecindad $(\delta_{S}(v)\leq p-1)$, por lo que $d_S(v)\leq p-1$, de esta forma, podemos deducir lo siguiente:

$$(p-1)|D_{k}^*|\geq \sum_{v\in D_{k}^*}d_v \hspace{3cm} (2)$$

%Aristas de D_k^* a S
%Hasta ahora son vecindades cerradas
Ahora, Sea $c_{u}'=|D_{k}^* \cap N_{G}(u)|$ para toda $u$ en $S$. Si $u\in D_{k}^* \cap S$, $c_u=c_{u}'+1$ y si $u\in S \setminus D_{k}^*, \hspace{.3cm} c_u=c_{u}'$, entonces

$$\sum_{u\in S}c_u = \sum_{u\in S\cap D_{k}^*} c_u' + 1 + \sum_{S\setminus D_k^*} c_u' \hspace{3cm}$$

$$\sum_{u\in S}c_u=|D_{k}^*\cap S| + \sum_{u\in S}c_{u}' \hspace{3cm} (3)$$

% conteo de aristas de S a d_k^*
Siguiendo el razonamiento anterior, sea $d_{v}'=|S\cap N_{G}(v)|$ para todo $v$ en $D_{k}^*$. Si $v\in S\cap D_{k}^*$, $d_v=d_{v}'+1$ y si $v\in D_{k}^*\setminus S \hspace{.3cm} D_v=d_{v}'$, entonces 

$$\sum_{v\in D_{k}^*}d_v = \sum_{v\in S\cap D_{k}^*} d_v' + 1 + \sum_{D_k^*\setminus S} d_v' \hspace{3cm}$$

$$\sum_{v\in D_{k}^*}d_v=|D_{k}^*\cap S|+\sum_{v\in D_{k}^*}d_v' \hspace{3cm} (4)$$

Sumamos las desigualdades (1) y (2)\\
%Terminamos los conteos y hacemos algebra con las desigualdades
\begin{align*}
\sum_{u\in S}c_u + (p-1)|D_{k}^*| &\geq k|S|+\sum_{v\in D_{k}^*}d_v\\
\text{Sustituimos las igualdades (3) y (4)}\\
|D_{k}^*\cap S| + \sum_{u\in S}c_{u}'+(p-1)|D_{k}^*| &\geq k|S|+|D_{k}^*\cap S|+\sum_{v\in D_{k}^*}d_v'\\
\sum_{u\in S}c_{u}'+(p-1)|D_{k}^*| &\geq k|S|+\sum_{v\in D_{k}^*}d_v'\\
\end{align*}

Notamos que $\sum_{u\in S}c_{u}' = \sum_{v\in D_{k}^*}d_v'$, ya que la suma anterior es la cantidad de aristas entre los conjutos $S$ y $D_k^*$, $c_u'$ cuenta desde el extremo que se encuentra en $S$ y $d_v'$ cuenta desde el extremo en $D_k*$, por lo que:\\

\begin{comment}
Notemos que $\sum_{u\in S}c_{u}' = \sum_{v\in D_{k}^*}d_v'$, ya que estamos haciendo el conteo de aristas entre el conjunto $\alpha(G)$-independiente y el conjunto $\gamma_{\times k}(G)$-dominante, el primer término parte del extremo que se encuentra en el conjunto $S$ y el lado derecho de la igualdad parte desde el extremo que se encuentra en $D_{k}^*$, por lo que:\\
\end{comment}


\begin{align*}
(p-1)|D_{k}^*| &\geq k|S|\\
|D_{k}^*|&\geq \frac{k|S|}{(p-1)}\\
\end{align*}

Por lo tanto

\begin{align*}
\gamma_{\times k}(G)&\geq \frac{k\alpha(G)}{p-1}\\
\end{align*}

 Con lo que hemos terminado la prueba.\\


 Como consecuencia directa del lema anterior, consideramos $p=3$, siendo este el valor que debe tomar $p$ para tener una gráfica libre de garras y con $\delta(G)\geq k-1$, por lo que, es el mejor valor que podemos obtener para la cota inferior del número de $k$-tupla dominación para una gráfica libre de garras, de esta forma, la cota queda como:

$$\gamma_{\times k}(G)\geq \frac{k\alpha(G)}{2}$$

El siguiente resultado presenta una cota para el número de $k$-tupla dominación en términos del número de independencia, pero en el sentido opuesto a la desigualdad anterior.\\

%\newpage

\textbf{Proposición 3.2} Sea $k\geq 2$ un entero. Para cualquier gráfica libre de garras $G$ de orden $n$ y $\delta(G)\ge k+1$,

$$\gamma_{\times k}(G)\leq n-\alpha(G)$$

\textbf{Dem:}\\

Sea $I$ un conjunto $\alpha(G)$-conjunto. Vamos a probar que $W=V(G)\setminus I$ es un conjunto $k$-tupla dominante de $G$.
%W es el conjunto que deja fuera a todos los del conjunto independiente.
Primero notemos que $\delta_W(x)\geq \delta(G)>k \hspace{.3cm} \forall x \in I$. Sea $v\in W$ y supongamos que $\delta_W(v) \leq k-2$, entonces $\delta(v)\geq 3$ ($k+1 - 3 = k-2$ ), es decir, $v$ al estar en $W$, pueden existir 3 vértices en $N(v)\cap I$, pero esto es una contradicción a que $G$ es libre de garras, entonces $\delta_W(v)\leq k-1$, por lo que $W$ es un conjunto $k$-tupla dominante de $G$, por lo tanto

$$\gamma_{\times k}(G)\leq |W|=n-\alpha(G)$$

$\hspace{15.5cm}\blacksquare$\\

La desigualdad anterior es justa para algunos valores de $k$. Para $k=2$ la desigualdad es justa para la gráfica $G=K_1 + C_4$, es decir, la rueda de orden 5. 

\begin{center}
\begin{tikzpicture}[scale=3.5, every node/.style={draw, circle, inner sep=2pt, minimum size=5mm}
]
\node (A1)[fill=blue!20] at (0,1) {$A_1$};
\node (A2)[fill=red!20] at (0,0) {$A_2$};
\node (A3)[fill=blue!20] at (1,0) {$A_3$};
\node (A4)[fill=red!20] at (1,1) {$A_4$};
\node (C)[fill=blue!20] at (.5,.5) {C};

\foreach \i [evaluate=\i as \nexti using int(\i+1)] in {1,...,3}{
\draw (A\i) -- (A\nexti);
\draw (A\i) -- (C);-
}
\draw (A4) -- (A1);
\draw (A4) -- (C);
\end{tikzpicture}

\vspace{0.5cm}

Figura 10: \small{Rueda de orden 5, siendo $I=\{A_2,A_4\}$ y $D_k^*=\{A_1, A_3, C\}$.}


\end{center}
Podemos exhibir el conjunto $\gamma_{\times k}(G)$-dominante de $G$, lo denotaremos como $D_k^*=\{A_1, A_3, C\}$, así como el conjunto independiente mámimal, $I=\{A_2, A_4\}$. Así $\gamma_{\times k}(G)=3$, $\alpha(G)=2$ y $n=5$, entonces:
$$3\leq 5-2=3$$
%\begin{align*}
%5&\leq 5-2
%\end{align*}
%Sección de cotas que dejan de usar números de independencia y comienzan a usar grados, en particular, grados máximos.

\newpage

\section{Cotas para el número de k-tupla dominación en términos del grado de una gráfica}


En esta sección se establece una cota superior y una cota inferior al número de $k$-tupla dominación utilizando el grado máximo de la gráfica, $\Delta(G)$. A contiuación, se presenta la cota superior.\\

\textbf{Teorema 3.3} Sea $k\geq 2$ un entero. Para cualquier gráfica $G$ con $\delta(G)\geq k$,\\
$$\gamma_{\times k}(G)\leq \frac{k\Delta(G)n}{k\Delta(G)+1}$$
\textbf{Dem:}\\
Sea $D$ un $\gamma_{\times k}(G)$-conjunto. Consideramos los siguientes subconjuntos:\\

\begin{multicols}{3}
%Todos los vértices en D que tienen toda su vecindad en D
$D^*=\{x\in D:N(x)\subseteq D\}$ 
\columnbreak
\hspace{1 cm} \& \hspace{1 cm}
\columnbreak
%Existe al menos un vértice de su vecindad que está fuera de D
$D^==\{x\in D:\delta_D(x)=k-1\}$
\end{multicols}
%Intersección vacia de los dos conjuntos
%Si un vertice u \in N(x) también está en D^=
Notemos que, si $x\in D^*$, $\delta_{D^*}(x)=\delta_G(x)\geq \delta(G)\geq k \neq k-1$, entonces

$$D^*\cap D^= =\emptyset \hspace {1 cm} (1)$$

Por otro lado, sea $x \in D^*$, entonces $D\setminus \{x\}$ no es $k$-tupla dominante, por lo que existe un vértice $u \in N(x)$ tal que $\delta_D(u)=k-1$, de este modo $u \in D^=$ y 

$$N(x)\cap D^= \neq \emptyset \hspace {1cm} (2)$$


%Necesitamos agregar la definicón de (A, B)-aristas

%Podemos contar las aristas de un conjunto a otro conjunto
Con lo anterior, tenemos aristas de la forma $(x,y)$ tales que $x\in D^*$ y $y\in D^=$, de esta forma $\delta_D(y)=k-1$ y como $D^*\subset D$ entonces, a lo más $\delta_{D^*}(y)=k-1$. Al hacer el conteo de las $(D^*, D^=)$-aristas, es a lo más $(k-1)|D^=|$ y por $(2)$, por cada vértice de $D^*$ hay al menos una $(D^*,D^=)$-arista, por lo que $|D^*|\le (k-1)|D^=|$. Por $(1)$ tenemos que $D^= \subseteq D\setminus D^*$, así $|D^=| \le |D\setminus D^*|$, entonces:\\


%Conteo de aristas de D^= a D*
\begin{align*}
|D^*|&\leq (k-1)|D^=|\\
&\leq (k-1)(|D|-|D^*|)\\
& = (k-1)|D|-(k-1)|D^*|\\
\Longrightarrow k|D^*|&\leq (k-1)|D|\\
\Longrightarrow |D^*|&\leq \frac{(k-1)}{k}|D| & (*)
\end{align*}
%Los vértices fuera del conjunto k-tupla dominante, dominan a G-D* (todos los vétices menos los vértices que están dentro del conjunto dominante junto con toda su vecindad)


%Necesitamos agregar la definición de restar un conjunto de vértices a una gráfica

Por la definición del conjunto $D^*$, $V(G)\setminus D$ es un conjunto dominante de $G-D^*$. Por lo tanto: \\


$$\gamma(G-D^*)\le |V(G)\setminus D|$$


Por el teorema "2" de Cap. 0, tenemos que:\\

 $$\frac{n-|D^*|}{\Delta(G-D^*)+1}\leq \gamma(G-D^*) \le n - |D|$$

Además, por la definición del grado máximo de una subgráfica, tenemos que $\Delta(G-D^*) \le \Delta(G)$, entonces:\\

$$\frac{n-|D^*|}{\Delta(G)+1}\leq \frac{n-|D^*|}{\Delta(G-D^*)+1}$$



%Desigualdad 2 para acotar el numero de $k$-tupla dominación
\begin{align*}
&\Longrightarrow\frac{n-|D^*|}{\Delta(G)+1}\leq n-|D| \\
&\Longrightarrow n-|D^*|\leq (n-|D|)(\Delta(G)+1)\\
& \Longrightarrow n-|D^*| \leq n\Delta(G)+n-|D|(\Delta(G)+1)\\
& \Longrightarrow |D^*|-n \geq |D|(\Delta(G)+1)-n\Delta(G)-n\\
& \Longrightarrow |D^*| \geq |D|(\Delta(G)+1)-n\Delta(G) & (**)
\end{align*}

De (*) y (**) tenemos lo siguiente:\\
\begin{align*}
\frac{(k-1)}{k}|D|&\geq |D|(\Delta(G)+1)-n\Delta(G)\\
\Longrightarrow n\Delta(G)&\geq |D|(\Delta(G)+1)-\frac{(k-1)}{k}|D|\\
\Longrightarrow n\Delta(G)&\geq |D|(\Delta(G)+1)-\frac{(k-1)}{k})\\
\Longrightarrow n\Delta(G)&\geq |D|(\frac{k\Delta(G)+k-(k-1)}{k})\\
\Longrightarrow n\Delta(G)&\geq|D|(\frac{k\Delta(G)+1}{k})\\
\Longrightarrow \frac{n\Delta(G)k}{k\Delta(G)+1}&\geq |D| = \gamma_{\times k}(G)
\end{align*}


Lo que es la desigualdad buscada y termina la prueba.\\

%Referencia 5, sus definiciones y el teorema del que proviene
La desigualdad anterior, acota por arriba al número de $k$-tupla dominación y los siguientes dos resultados son cotas inferiores del mismo número. Comenzamos yendo a la \cite{harary2000double}, el cual solo funciona como una cita y antesala ante el "refinamieno" del último teorema.\\

%Definiciones necesarias para el resultado

%Doble dominante, caso más particular de la k-tupla dominación, k=2  /Usa vecindad uy grado cerrado

%\newpage

%Caso particular de k=2, doble dominación
\textcolor{Violet}{\textbf{Teo.6}} Si $G$ no tiene vértices aislados, entonces\\

$$\gamma_{\times 2}(G)\geq\frac{2n}{\Delta(G)+1}$$

\textbf{Dem:}\\
%La demostración es un doble conteo de aristas

%Popiedades de grados y del conjunto dominante, así como el incio que marca el conteo de aristas desde S
Sea $S$ un $\gamma_{\times 2}(G)$-conjunto de $G$ y sea $t$ el número de aristas en $G$ de la forma $(v,u)$, tales que $v\in S$ y $u\in V(G)\setminus S$. Como $S$ es $\gamma_{\times 2}(G)$-conjunto, $\forall v\in S, \hspace{2mm} |N(v)\cap S|\geq 1$ y como $\Delta(G)\geq \delta(v)$ tenemos $\Delta(G)-1\geq \delta_{V(G)\setminus S}(v)$. Esto sucede para cada vértice dentro de $S$, por lo que:\\

%Conteo de aristas acotado desde S
$$t\leq (\Delta(G)-1)|S|=(\Delta(G)-1)\gamma_{\times 2}(G) \hspace{2cm} (1)$$

%Comenzamos el conteo desde V(G)\S
Por otro lado, $\forall u\in V(G)\setminus S, \hspace{2mm} |N(u)\cap S|\geq 2$. Considerando la cantidad de aristas $(u,v)$ antes mencionadas, tenemos que:\\

%Cota inferior desde V(G)\S
$$t\geq 2|V(G)\setminus S|=2(n-\gamma_{\times 2}(G)) \hspace{2.7cm} (2)$$

De $(1)$ y $(2)$ tenemos\\
\begin{align*}
2(n-\gamma_{\times 2}(G)) & \le(\Delta(G)-1)\gamma_{\times 2}(G)\\
2n-2 \gamma_{\times 2}(G) & \le \gamma_{\times 2}(G)\Delta(G) - \gamma_{\times 2}(G)\\
2n & \le \gamma_{\times 2}(G)\Delta(G) + \gamma_{\times 2}(G)\\
2n & \le \gamma_{\times 2}(G)(\Delta(G) + 1)\\
\Longrightarrow \frac{2n}{\Delta(G)+1}&\leq \gamma_{\times 2}(G)
\end{align*}

Este resultado, es un caso particular para la $k$-tupla dominación cuando $k=2$. Se puede generalizar el resultado anterior para $k\geq 2$.\\

%Caso general que falta por demostrar.
\textcolor{Violet}{\textbf{Teo.14}} Sea $G$ una gráfica tal que $\delta(G)\geq k-1$, entonces\\
$$\gamma_{\times k}(G)\geq \frac{kn}{\Delta(G)+1}$$


El siguiente y último resultado, presenta una mejora parcial al anterior, pues solo es para gráficas que cumplan que $(\Delta(G)+1)|V_\Delta(G)|<kn$, donde $V_\Delta(G)\{v\in V(G): \delta_{V(G)}(v)=\Delta(G)\}$\\

%\newpage

\textbf{Teorema 3.4} Sea $k\geq 2$ un entero. Sea $G$ una gráfica y sea $V_\Delta(G)=\{v\in V(G): \delta_{V(G)}(v)=\Delta(G)\}$. Si $\delta(G)\geq k-1$, entonces\\

$$\gamma_{\times k}(G)\geq \frac{kn-|V_\Delta(G)|}{\Delta(G)}$$

\textbf{Dem:}\\

%Separación de conjunto D y los vértices fuera de él.
Dados dos conjuntos $D_1, D_2 \subseteq V(G)$, sea $E(D_1,D_2)=\{uv\in E(G): u\in D_1 \hspace{2mm} v\in D_2\}$. Sea $D$ un $\gamma(G)$-conjunto, por lo que $\forall v\in (V(G)\setminus D)$ es dominado por al menos $k$ vértices de $D$ ($\delta_D(v)\geq k$), por lo que $k|(n-|D|)|$ es la cantidad mínima de aristas que podemos tener entre $D_1$ y $D_2$, entonces

%Acotamiento del conteo de aristas desde fuera del conjunto dominante
$$k|(n-|D|)|\leq |E(V(G)\setminus D,D)|$$

%Conteo desde dentro del conjunto dominante
Ahora, para cada $v\in D$, $\delta(v)-\delta_D(v)$, son todos sus vecinos que tiene fuera de $D$. La suma de lo anterior para todos los vértices de $D$, es equivalente a contar las aristas que existen entre $V(G)\setminus D$ y $D$, por lo que:

%Relación de la primer y desigualdad con la segunda obtenida
\begin{align*}
k(n-|D|)&\leq |E(V(G)\setminus D),D)|\\
&\leq \sum_{v\in D}(\delta(v)-\delta_D(v)) \hspace{3cm} (1)
\end{align*}

%Grado de vértices con grado máximo y en el conjunto dominante
Notemos que sucede para los vértices que son de grado máximo y están en $D$. Para cada $v\in D\cap V_\Delta(G)$, tienen grado máximo y tienen al menos $k-1$ vecinos dentro del conjunto $D$, $\delta(v)-\delta_D(v)\leq \Delta(G)-(k-1)=\Delta(G)-k+1 \hspace{1cm} *$\\

%Vértices en D que no son de grado máximo
Para todo $v\in (D\setminus V_\Delta(G))$, $\delta(v)\leq \Delta(G)-1$ y como $D$ es $k$-tupla dominante, $\delta_D(v)\geq k-1 \hspace{1cm} **$\\
 entonces\\

%Suma de las últimas dos desigualdades obtenidas
\begin{align*}
\delta(v)+k-1 &\leq \Delta(G)-1+\delta_D(v)\\
\delta(v)-\delta_D(v) &\leq \Delta(G)-k
\end{align*}
Sumando $*$ y $**$, además tomando la cardinalidad de donde se encuentra $v$, tenemos lo siguiente:\\
\begin{align*}
\sum_{v\in D}\delta(v)-\delta_D(v) \leq& |D\cap V_\Delta(G)|(\Delta(G)-k+1) + (|D|-|D\cap V_\Delta(G)|)(\Delta(G)-k)\\
\text{Usando (1)}\\
\Longrightarrow k(n-|D|) \leq& |D\cap V_\Delta(G)|(\Delta(G)-k+1) + (|D|-|D\cap V_\Delta(G)|)(\Delta(G)-k)\\
\Longrightarrow kn - k|D| \leq& |D\cap V_\Delta(G)|\Delta(G) - |D\cap V_\Delta(G)|k + |D\cap V_\Delta(G)| + |D|\Delta(G)\\
& - |D|k - |D\cap V_\Delta(G)|\Delta(G) - |D\cap V_\Delta(G)|k\\
\Longrightarrow kn \leq& |D\cap V_\Delta(G)| + |D|\Delta(G)\\
\Longrightarrow kn - |D\cap V_\Delta(G)| \leq& |D|\Delta(G)\\
\Longrightarrow \frac{kn - |D\cap V_\Delta(G)|}{\Delta(G)} \leq& |D|\\
\end{align*}
Como $|D|=\gamma_{\times k}(G)$ y $|D\cap V_\Delta(G)|\leq |V_\Delta(G)|$, podemos concluir que:\\
$$\frac{kn - |V_\Delta(G)|}{\Delta(G)} \leq \gamma_{\times k}(G)$$\\

Es la desigualdad buscada, por lo que hermos terminado la demostración.\\

Por último, mostramos un ejemplo en el que podemos ver el caso donde la desigualdad es justa para la gráfica $G$ y $k=2$.\\


\begin{center}
\begin{tikzpicture}[
    scale=1.5,
    every node/.style={
        circle,
        draw,
        inner sep=1.2pt,
        minimum size=4.2mm
    }
]

\foreach \i/\x in {1/0, 2/1.6, 3/3.2, 4/4.8, 5/6.4} {
   
    \node[fill=white] (v\i) at (\x,0) {$v_\i$};
}
\foreach \i/\y in {7/1.6,6/3.2}{
	
	\node[fill=white] (v\i) at (3.2,\y) {$v_\i$};

}

\foreach \i in {1,...,5} {
	\pgfmathtruncatemacro{\next}{\i+1}
	\draw (v\i) -- (v\next);
}

\draw (v6) -- (v1);
\draw (v2) -- (v7);
\draw (v4) -- (v7);
\end{tikzpicture}\\
\end{center}

%\newpage

Para $k=2$, podemos ver el conjunto $k$-tupla dominante $D=\{v_1, v_2, v_4, v_5\}$. Por lo que podemos determinar los valores restantes para poder sustituir sobre la desigualdad del teorema anterior, $k=2$, $|V(G)|=7$, $\Delta(G)=3$, $|V_\Delta(G)|=2$, $\gamma_{\times 2}(G)=4$. Entonces:\\

\begin{align*}
\frac{2(7)- 2}{3}&\leq 4\\
\frac{14- 2}{3}&\leq 4\\
\frac{12}{3}&\leq 4\\
4 & = 4\\
\end{align*}


Por último, podemos notar que la desigualdad es justa para cualquier gráfica $G=K_k + H$, obtenidad de la gráfica completa $K_k$ y cualquier gráfica $H$ de orden $|V(H)|=k>2$ y $\Delta(H)<k-1$. Pra esta gráfica $G$, tenemos que $\gamma_{\times k}(G)=|V_\Delta(G)|=k$, $|V(G)|=2k$ y $\Delta(G)=2k-1$.


\newpage


\begin{thebibliography}{9}
        
		\bibitem{chartrand2012first} 
        G. Chartrand and P. Zhang, \textit{A First Course in Graph Theory}. Dover Publications, 2012.

		\bibitem{chartrand2015graphs} 
        G. Chartrand, L. Lesniak, and P. Zhang, \textit{Graphs \& Digraphs} (6th ed.). CRC Press, 2015.

        \bibitem{cabrera2022note} %8
        A. Cabrera-Martínez, \textit{A note on the $k$-tuple domination number of graphs}. Ars Math. Contemp. 22 (2022) P4.03.
        
        \bibitem{cabrera2022some} %principal
        A. Cabrera-Martínez, \textit{Some new results on the $k$-tuple domination number of graphs}. RAIRO-Oper. Res. 56 (2022) 3491-3497.
        
        \bibitem{hansberg2020multiple} %6
    	A. Hansberg and L. Volkmann, \textit{Multiple domination}, in Topics in Domination in Graphs. Developments in Mathematics. Springer (2020) 151-203.

    	\bibitem{favaron2001domination} %15
    	O. Favaron, M.A. Henning, J. Puech and D. Rautenbach, \textit{On domination and annihilation in graphs with claw-free blocks}. Discrete Math. 231 (2001) 143-151.

    	\bibitem{klasing2004hardness} %16
    	R. Klasing and C. Laforest, \textit{Hardness results and approximation algorithms of k-tuple domination in graphs}. Inform. Process. Lett. 89 (2004) 75-83.

    	\bibitem{Haynes2020topics} %1
    	T.W. Haynes, S.T. Hedetniemi and M.A. Henning,Topics in domination in graphs, in \textit{Developments in Mathematics. Springer (2020).}


    	\bibitem{harary2000double} %5
    	F. Harary and T.W. Haynes, \textit{Double domination in graphs}. Ars Combin. 55 (2000) 201-213.

    	\bibitem{alipour2020upper} 
        S. Alipour, A. Jafari, and M. Saghafian, \textit{Upper bounds for k-tuple (total) domination numbers of regular graphs}. Bulletin of the Iranian Mathematical Society 46 (2020) 573-577.
    \end{thebibliography}

	








	


	
	
	
\end{document}
